% scuc_manual.tex —— MIPSolvers SCUC 模块技术手册（工程参考级）
% 规范：docs/modules/README.md 与 docs/modules/THEORY_WRITING_GUIDE.md。
% 编译（本目录下，XeLaTeX 两遍）：
%   xelatex -interaction=nonstopmode scuc_manual.tex
%   xelatex -interaction=nonstopmode scuc_manual.tex
\documentclass{ctexart}
\usepackage{../../_manual_common/mipsolvers_manual}

\title{MIPSolvers SCUC 模块技术手册}
\author{MIPSolvers 文档组}
\date{2026 年 9 月 13 日}

\begin{document}
\maketitle

\begin{abstract}
本手册是 MIPSolvers 安全约束机组组合（Security-Constrained Unit Commitment, SCUC）
模块的工程参考级技术文档。模块源码位于 \srcpath{src/scuc/}，公共头文件位于
\srcpath{include/mipsolvers/scuc/}；它以 JSON 为交换格式，实现《南方区域电力市场
现货交易规则（2025 V1.0）》§2.6 的日前出清三阶段流水线：SCUC（MILP）→
SCED（固定组合的 LP）→ LMP（节点边际电价，由 SCED 对偶提取），并附带合成算例
构造器（case\_builder）与两个命令行工具（\texttt{scuc\_solve}、
\texttt{scuc\_case\_builder}）。

手册分两层：第 3 章为通用理论层（标准模型、MILP 松弛、LMP 原理，非代码转写）；
第 4 章为源码等价转写层（JSON→MILP 的每个约束族、系数与索引集均按当前源码
如实转写）；第 5 章给出数值验证证据，第 6 章列出 JSON 接口字段，第 7 章声明
验证状态与边界局限，第 8 章为跨模块共享的工业级评价基线。所有
\srcpath{文件:函数名} 锚点均为仓库相对路径、禁止行号，由
\texttt{tools/doc\_anchor\_check.py} 回归校验。
\end{abstract}

\tableofcontents
\newpage

\section{阅读指南与约定}

\subsection{数据来源与诚实口径}

本手册实现层（第 4 章）的唯一事实来源是下列源码文件，转写前已逐函数通读；
凡代码未消费的输入字段、回退路径、硬编码阈值均在相邻正文明确声明：

\begin{center}
\begin{footnotesize}
\begin{tabular}{@{}ll@{}}
\toprule
\textbf{文件} & \textbf{内容}\\
\midrule
\srcpath{include/mipsolvers/scuc/scuc.hpp} & 数据结构、配置与公共 API 声明\\
\srcpath{include/mipsolvers/scuc/case\_builder.hpp} & 算例构造器声明\\
\srcpath{src/scuc/scuc.cpp} & JSON 解析、MILP 装配、三阶段求解、结果序列化\\
\srcpath{src/scuc/case\_builder.cpp} & 合成算例（3bus/6bus/ieee39/ieee118）生成\\
\srcpath{src/scuc/main.cpp} & \texttt{scuc\_solve} 命令行入口\\
\srcpath{src/scuc/case\_builder\_main.cpp} & \texttt{scuc\_case\_builder} 命令行入口\\
\bottomrule
\end{tabular}
\end{footnotesize}
\end{center}

诚实口径约定：
\begin{itemize}
  \item 公式按\textbf{代码实际执行}转写；代码与《南方区域电力市场现货交易规则
        （2025 V1.0）》§2.6 条文存在简化差异之处，以"部分实现"或 gapnote 标注，
        不暗示代码即规则完整模型；
  \item 数值结论只引用第 5 章所列真实测试目标与可复现命令；
  \item 锚点一律为 \texttt{文件:函数名}（匿名命名空间内函数以源文件名加函数名
        引用，lambda 记作 \texttt{文件名:lambda名（lambda）}），禁止行号。
\end{itemize}

\subsection{索引空间与单位制}

\begin{itemize}
  \item \textbf{索引}：母线、支路、机组、负荷、风电、光伏、储能、直流线、断面
        一律为 \textbf{0 起始}整数索引；\texttt{num\_buses} 缺省时由所有元件
        母线编号最大值加 1 推断（\srcpath{src/scuc/scuc.cpp:scuc\_from\_json}）。
  \item \textbf{时间}：调度时段数 $T=$ \fld{num\_periods}，时段长
        $\Delta t=$ \fld{period\_length\_hr}（小时）；每小时调度区间数
        $\mathrm{iph}=\mathrm{round}(1/\Delta t)\ge 1$，组合时段数
        $T_c=\max(1,\lfloor T/\mathrm{iph}\rfloor)$。启停类二进制变量按组合
        时段索引 $h=0,\dots,T_c-1$，出力类变量按调度时段 $t=0,\dots,T-1$，
        二者经 $h=\lfloor t/\mathrm{iph}\rfloor$ 耦合
        （\srcpath{src/scuc/scuc.cpp:setup\_var\_index}）。
  \item \textbf{单位}：功率 MW，能量 MWh，电抗 pu（100 MVA 基准口径由算例
        数据自带），价格与成本以货币单位（记 \$）/MWh、\$/h 或 \$/次 计；
        爬坡率为 MW/min；最小启停时间为小时。
  \item \textbf{松弛母线}：直流潮流灵敏度（PTDF）计算固定取母线 0 为平衡
        母线（\srcpath{src/scuc/scuc.cpp:compute\_ptdf}）。
\end{itemize}

\subsection{排版宏}

参数表使用 \texttt{paramtable} 环境（字段名|符号|单位|典型范围|默认值|说明）；
字段名以 \verb|\fld{}| 排版；源码锚点以 \verb|\srcpath{}| 排版；结构体元信息以
\verb|\compmeta{名称}{头文件}| 排版；理论章三要素（theorynote、与实现的对应表、
gapnote）按 \texttt{docs/modules/THEORY\_WRITING\_GUIDE.md} 执行。

\section{功能定位与架构}

\subsection{功能定位}

scuc 模块是面向电力市场日前出清的\textbf{应用层}模块：它把一份描述电网、机组、
负荷、可再生、储能与市场参数的 JSON 文档装配成一个混合整数线性规划（MILP），
调用 \texttt{mipsolvers::engine} 注册的任一 MILP 后端求解，再依次执行固定组合
的经济调度（SCED，LP）与节点电价（LMP）计算，最终把三阶段结果序列化为 JSON。
模块自身不实现任何 LP/MILP 算法；全部求解能力来自 engine 模块（见 engine 手册）。

\subsection{入口一览}

\begin{footnotesize}
\begin{longtable}{@{}L{3.4cm}L{4.6cm}L{6.2cm}@{}}
\toprule
\textbf{入口} & \textbf{锚点} & \textbf{说明}\\
\midrule
\endhead
CLI \texttt{scuc\_solve} & \srcpath{src/scuc/main.cpp:main} &
读入 JSON → 三阶段求解 → 写出 JSON；支持 \texttt{--solver}、\texttt{--no-sced}、
\texttt{--no-lmp}、\texttt{--indent}；SCUC 未收敛时返回非零退出码\\
CLI \texttt{scuc\_case\_builder} & \srcpath{src/scuc/case\_builder\_main.cpp:main} &
生成内置算例 JSON；支持 \texttt{--case 3bus|6bus|ieee39|ieee118}、\texttt{--T}、
\texttt{--dt}、\texttt{--wind}、\texttt{--solar}、\texttt{--storage}、
\texttt{--output}、\texttt{--compact}\\
库 API：解析 & \srcpath{src/scuc/scuc.cpp:scuc\_from\_json} &
JSON 字符串 → \texttt{SCUCInput}；语法错误抛 \texttt{json::parse\_error}，
类型不匹配抛 \texttt{json::type\_error}\\
库 API：求解 & \srcpath{src/scuc/scuc.cpp:scuc\_solve} &
\texttt{SCUCInput} → 三阶段流水线 → \texttt{SCUCOutput}\\
库 API：序列化 & \srcpath{src/scuc/scuc.cpp:scuc\_output\_to\_json} &
\texttt{SCUCOutput} → JSON 字符串（\texttt{indent=0} 紧凑、\texttt{2} 可读）\\
库 API：裸模型 & \srcpath{src/scuc/scuc.cpp:build\_scuc\_mip} &
只装配不求解，返回 \texttt{engine::MIPModel}，供基准程序注入自定义
\texttt{BCOptions}\\
算例构造 & \srcpath{src/scuc/case\_builder.cpp:build\_3bus\_case} 等 &
\texttt{build\_3bus\_case}/\texttt{build\_6bus\_case}/\texttt{build\_ieee39\_case}/\texttt{build\_ieee118\_case} 返回预填 \texttt{SCUCInput}；
\texttt{scuc\_input\_to\_json} 导出 JSON\\
\bottomrule
\end{longtable}
\end{footnotesize}

构建开关：\texttt{MIPSOLVERS\_BUILD\_SCUC}（默认 ON）控制 \texttt{scuc\_solve}
与相关测试，\texttt{MIPSOLVERS\_BUILD\_SCUC\_CASE\_BUILDER}（默认 ON）控制
\texttt{scuc\_case\_builder}（根 \texttt{CMakeLists.txt}）。

\subsection{关键数据结构}

\compmeta{SCUCInput}{include/mipsolvers/scuc/scuc.hpp}
完整问题输入：\texttt{config}（求解配置）、\texttt{num\_buses} 与七个元件数组
（\texttt{generators/branches/loads/wind/solar/storage/dc\_lines}）、
\texttt{generator\_groups}、\texttt{sections}、\texttt{profiles}、
\texttt{initial\_status}。

\compmeta{SCUCConfig}{include/mipsolvers/scuc/scuc.hpp}
求解配置：求解器选择与回退、时段结构、MIP 间隙与时限、备用需求系数、
惩罚价格（VOLL/VOCC/M1/M2）、输电费、市场割平面与原始修复开关、
SCED/LMP 阶段开关与 LMP 邻域因子。

\compmeta{Generator}{include/mipsolvers/scuc/scuc.hpp}
可调度机组：出力上下限、爬坡率、最小启停时间、must-run、启停次数上限、
分段报价 \texttt{bid\_segments}、启动成本（含热/温/冷三态字段，当前实现只用
\texttt{startup\_cost}，见第 4 章声明）、空载成本、备用报价、一次调频系数
\texttt{pfr\_alpha}、群标识 \texttt{group\_id}。

\compmeta{Branch}{include/mipsolvers/scuc/scuc.hpp}
交流支路：\texttt{from/to}（0 起始）、电抗（pu）、热稳定限额（MW）、
\texttt{in\_service}。

\compmeta{StorageUnit}{include/mipsolvers/scuc/scuc.hpp}
储能/抽蓄：充放功率上下限、能量容量、效率（单效率或充/放分效率）、
SOC 初值/下限/末值要求、充放电报价、日循环上限、二进制模式指示开关。

\compmeta{Section}{include/mipsolvers/scuc/scuc.hpp}
监视断面：成员线路与权重的线性组合构成断面 PTDF，正/反向限额分别给定。

\compmeta{SCUCSolveResult}{include/mipsolvers/scuc/scuc.hpp}
单阶段结果：收敛标志、目标值、求解器名、用时、MIP 间隙、割平面数，
以及组合/启停、出力、备用、分段出力、可再生出力与弃电、储能充放与 SOC、
线路与断面潮流、系统级切负荷/弃电、成本分解等矩阵。

\compmeta{SCUCLMPResult}{include/mipsolvers/scuc/scuc.hpp}
LMP 结果：节点电价及其能量/阻塞分量 \texttt{[nb][T]}，时空汇总统计
（avg/max/min）。

\compmeta{SCUCOutput}{include/mipsolvers/scuc/scuc.hpp}
流水线总输出：\texttt{scuc}（MILP）、\texttt{sced}（LP，仅当
\fld{solve\_sced}）、\texttt{lmp}（仅当 \fld{solve\_lmp}）。

\subsection{三阶段流水线}

\begin{enumerate}
  \item \textbf{SCUC（MILP）}：\srcpath{src/scuc/scuc.cpp:solve\_scuc\_stage}
        → 装配（\srcpath{src/scuc/scuc.cpp:build\_formulation}）→
        MILP 求解（\srcpath{src/scuc/scuc.cpp:run\_milp}）→
        结果提取（\srcpath{src/scuc/scuc.cpp:extract\_result}）；
  \item \textbf{SCED（LP）}：\srcpath{src/scuc/scuc.cpp:solve\_sced\_stage}
        重新装配并把全部启停/组合二进制固定为 SCUC 解（四舍五入），作为 LP 重解；
        仅当 \fld{solve\_sced}\,=\,true 且 SCUC 收敛；
  \item \textbf{LMP}：\srcpath{src/scuc/scuc.cpp:solve\_lmp\_stage}
        在 SCED（或 SCUC，若 SCED 未运行）组合之上加 $\delta$ 邻域出力限幅
        重解 LP 并提取约束对偶；仅当 \fld{solve\_lmp}\,=\,true 且参考阶段收敛。
\end{enumerate}

% theory_unit_commitment.tex —— scuc 通用理论：安全约束机组组合与节点电价
% 本文件为理论层章节，遵循 docs/modules/THEORY_WRITING_GUIDE.md。

\section{机组组合理论：标准模型、MILP 松弛与节点电价}

\begin{theorynote}
本章为通用数学理论，按电力市场出清与机组组合（Unit Commitment, UC）领域的标准
模型体系撰写：SCUC 的标准 MILP 提法、分段线性化、直流潮流与 PTDF、备用与储能
约束族、最小启停时间的紧 formulations、以及 SCED/LMP 定价原理。覆盖范围对应
《南方区域电力市场现货交易规则（2025 V1.0）》§2.6 的日前出清模型框架
\cite{southern2025}。本章公式不要求逐式对应代码；与平台实现
（\texttt{src/scuc/}）的对应关系见本章末"与实现的对应"表，超出实现的内容以
"未实现扩展说明"框就地标记。
\end{theorynote}

\subsection{记号与约定}

\begin{center}
\begin{footnotesize}
\begin{longtable}{@{}ll@{}}
\toprule
\textbf{符号} & \textbf{含义}\\
\midrule
\endhead
$\mathcal G,\ \mathcal B,\ \mathcal L$ & 机组、母线、交流支路集合；$n_g,n_b,n_l$ 为其基数\\
$\mathcal W,\mathcal S,\mathcal H,\mathcal D^{dc}$ & 风电、光伏、储能、直流线集合\\
$\mathcal T=\{0,\dots,T-1\}$ & 调度时段；$\Delta t$ 为时段长（h）\\
$\mathcal T_c=\{0,\dots,T_c-1\}$ & 组合时段；$h(t)=\lfloor t/\mathrm{iph}\rfloor$ 为耦合映射\\
$u_{g,h},y_{g,h},z_{g,h}\in\{0,1\}$ & 组合（在网）、启动、停机二进制变量\\
$P_{g,t}\ge 0$ & 机组 $g$ 在时段 $t$ 的总出力（MW）\\
$p_{g,k,t}\ge 0$ & 第 $k$ 报价段出力（MW），$k=1,\dots,K$\\
$R^{spin}_{g,t},R^{up}_{g,t},R^{dn}_{g,t}\ge 0$ & 旋转备用、上调、下调调节备用（MW）\\
$\underline P_g,\bar P_g$ & 最小/最大技术出力（MW）\\
$R^+_g,R^-_g$ & 上/下爬坡率（MW/min）\\
$T^U_g,T^D_g$ & 最小开机/停机时间（h，取整后以组合时段计）\\
$C^U_g,C^{NL}_g$ & 启动成本（\$/次）、空载成本（\$/h）\\
$c_{g,k}$ & 第 $k$ 段边际报价（\$/MWh）\\
$D_{b,t}$ & 母线 $b$ 在时段 $t$ 的负荷（MW）\\
$x_l,\bar F_l$ & 支路电抗（pu）与热稳定限额（MW）\\
$\mathrm{PTDF}_{l,b}$ & 支路 $l$ 对母线 $b$ 注入的功率传输分布因子\\
$M_1,M_2$ & 线路/断面松弛罚系数、可再生弃电罚系数（\$/MW、\$/MWh）\\
$c^{shed},c^{curt}$ & 失负荷价值 VOLL、过剩惩罚 VOCC（\$/MWh）\\
\bottomrule
\end{longtable}
\end{footnotesize}
\end{center}

全章未特别说明时，求和指标 $g\in\mathcal G$、$t\in\mathcal T$、$h\in\mathcal T_c$。

\subsection{SCUC 标准模型}

\subsubsection{问题提法}

\begin{definition}[SCUC]
安全约束机组组合问题是在调度时界 $\mathcal T$ 上同时决定各机组的启停状态
$u_{g,h}$ 与出力计划 $P_{g,t}$，最小化总运行成本，满足：（i）系统功率平衡与
备用需求；（ii）机组物理约束（出力限、爬坡、最小启停、启停次数）；（iii）储能
与直流线动态约束；（iv）输电网络安全约束（直流潮流口径的线路与断面限额）。
由于含二进制变量，SCUC 是混合整数线性规划（MILP），属 NP-hard
\cite{wood2013}。
\end{definition}

\subsubsection{目标函数}

标准日前出清目标为总成本最小化：
\begin{equation}
\min\ \underbrace{\sum_{g,t}\Big(C^{NL}_g\,u_{g,h(t)}\,\tfrac{1}{\mathrm{iph}}
  + \sum_{k} c_{g,k}\,p_{g,k,t}\,\Delta t\Big)}_{\text{空载 + 分段能量成本}}
+ \underbrace{\sum_{g,h} C^U_g\, y_{g,h}}_{\text{启动成本}}
+ \underbrace{\sum_{g,t}\big(c^{spin}_g R^{spin}_{g,t}+c^{up}_g R^{up}_{g,t}
  +c^{dn}_g R^{dn}_{g,t}\big)\Delta t}_{\text{备用成本}}
+ \Pi,
\label{eq:scuc-obj}
\end{equation}
其中 $\Pi$ 汇总各类惩罚项：切负荷 $c^{shed}\,\Delta t\,P^{shed}_t$、
过剩 $c^{curt}\,\Delta t\,P^{curt}_t$、线路/断面潮流松弛 $M_1(\sigma^+_{l,t}
+\sigma^-_{l,t}+\cdots)$、可再生弃电 $M_2\,\Delta t\,(P^{wc}_{w,t}+P^{sc}_{s,t})$，
以及可选的储能充放电报价项与输电费项。惩罚系数量级需显著高于正常能量价格
（典型 VOLL $\sim 10^4$ \$/MWh），使松弛仅在物理不可行时被使用；该设计即
"软约束 + 大 M 罚"的市场出清惯例 \cite{southern2025}。

\begin{remark}[惩罚口径的诚实声明]
大 M 罚使模型在物理不可行时仍返回"带松弛的最优解"。判读结果时必须先检查
松弛量（切负荷、弃电、线路松弛）是否为零量级，再接受调度方案；本模块在输出
JSON 中单独给出 \texttt{load\_shedding}、\texttt{gen\_curtailment} 与成本分解
中的 \texttt{penalty} 项用于该检查。
\end{remark}

\subsubsection{功率平衡与备用}

系统功率平衡（单区域、损耗忽略）：
\begin{equation}
\sum_g P_{g,t}+\sum_w P_{w,t}+\sum_s P_{s,t}
+\sum_{s'}\big(P^{dis}_{s',t}-P^{ch}_{s',t}\big)+\sum_d P^{dc}_{d,t}
+P^{shed}_t-P^{curt}_t=\sum_b D_{b,t},\quad \forall t.
\label{eq:balance}
\end{equation}

备用需求按系统负荷比例给定（§2.6.3.2--2.6.3.4 口径）：
\begin{align}
\sum_g R^{spin}_{g,t}&\ge r^{spin}\,D_t, &
\sum_g R^{up}_{g,t}&\ge r^{up}\,D_t, &
\sum_g R^{dn}_{g,t}&\ge r^{dn}\,D_t, \label{eq:reserve}\\
\sum_g \big(P_{g,t}-\underline P_g u_{g,h(t)}\big)&\le D_t - r^{neg}D_t, &
\sum_g \alpha^{pfr}_g \bar P_g u_{g,h(t)}&\ge R^{pfr}, && \nonumber
\end{align}
其中 $D_t=\sum_b D_{b,t}$。负备用约束的含义是：在线机组压到最小技术出力后的
总下调空间不得侵入负备用需求；一次调频（PFR）备用按系数
$\alpha^{pfr}_g\bar P_g$ 计入在线容量。

\subsubsection{机组约束族}

\paragraph{出力耦合与分段。}
\begin{equation}
P_{g,t}=\underline P_g u_{g,h(t)}+\sum_{k=1}^{K} p_{g,k,t},\qquad
\underline P_g u_{g,h(t)}\le P_{g,t}\le \bar P_g u_{g,h(t)},\qquad
0\le p_{g,k,t}\le q_{g,k}\,u_{g,h(t)},
\label{eq:gen-coupling}
\end{equation}
其中 $q_{g,k}$ 为第 $k$ 段容量。该"段变量 + 等式聚合"形式把分段线性成本
写成显式线性结构，见 \ref{sec:pwl} 节。

\paragraph{启停逻辑。}
\begin{equation}
u_{g,h}-u_{g,h-1}=y_{g,h}-z_{g,h},\qquad y_{g,h}+z_{g,h}\le 1,
\label{eq:logic}
\end{equation}
首时段以初始在网状态 $u_{g,0^-}$ 替换 $u_{g,-1}$；must-run 机组固定
$u_{g,h}=1$。

\paragraph{最小启停时间。}
记 $TU_g=\lceil T^U_g\rceil$、$TD_g=\lceil T^D_g\rceil$（组合时段数），
"启动触发在线"形式为 \cite{ostrowski2012}：
\begin{equation}
\sum_{\tau=\max(0,h-TU_g+1)}^{h} y_{g,\tau}\le u_{g,h},\qquad
\sum_{\tau=\max(0,h-TD_g+1)}^{h} z_{g,\tau}\le 1-u_{g,h},\qquad \forall h.
\label{eq:minupdown}
\end{equation}
时界起点处求和被截断（$\tau\ge 0$），初始状态历史不进入约束。

\begin{remark}[等价紧形式]
文献中常见的另一族 $\sum_{\tau=h}^{h+TU_g-1} z_{g,\tau}\le 1-u_{g,h}$（前视
形式）与 \eqref{eq:minupdown} 在整数意义下等价，LP 松弛强度不同；
Rajan--Takriti 证明最小启停 polytope 存在无附加变量的完全描述
\cite{rajan2005}，Ostrowski 等给出各类形式的 LP 收紧比较
\cite{ostrowski2012}，Morales-España 等给出"紧且紧凑"的三二进制变量形式
\cite{morales2013}。当前实现采用 \eqref{eq:minupdown} 的后视形式。
\end{remark}

\paragraph{爬坡。}
\begin{equation}
P_{g,t}-P_{g,t-1}\le 60\Delta t\,R^+_g + \bar P_g\, y_{g,h(t)},\qquad
P_{g,t-1}-P_{g,t}\le 60\Delta t\,R^-_g + \bar P_g\, z_{g,h(t)},
\label{eq:ramp}
\end{equation}
首时段以初始出力 $P_{g,0^-}$ 替换 $P_{g,-1}$。右端的
$\bar P_g y$（$\bar P_g z$）项是启动（停机）时段放开爬坡限制的 big-M 处理：
启动时段机组可从 0 爬至任意 $\le\bar P_g$ 的出力。更紧的形式用
$\min(\underline P_g+60\Delta t R^+_g,\bar P_g)$ 替代 $\bar P_g$（见
\ref{sec:tightening} 节）。

\paragraph{备用耦合。}
\begin{align}
P_{g,t}+R^{spin}_{g,t}+R^{up}_{g,t}&\le \bar P_g u_{g,h(t)},&
R^{dn}_{g,t}&\le P_{g,t}-\underline P_g u_{g,h(t)},\nonumber\\
R^{spin}_{g,t}&\le (\bar P_g-\underline P_g)u_{g,h(t)},&
R^{up}_{g,t}&\le 10 R^+_g\, u_{g,h(t)},\quad
R^{dn}_{g,t}\le 10 R^-_g\, u_{g,h(t)}.
\label{eq:reserve-coupling}
\end{align}
末两式对应"10 分钟可调"口径：调节备用上限取 10 分钟爬坡能力（
$10 R^{\pm}_g$，MW），这是§2.6 类规则中调节备用响应时间的标准近似
\cite{southern2025}。

\subsubsection{储能约束族}

\begin{align}
E_{s,t}&=E_{s,t-1}+\eta^{ch}_s \Delta t\, P^{ch}_{s,t}
-\frac{\Delta t}{\eta^{dis}_s} P^{dis}_{s,t},&
\underline E_s&\le E_{s,t}\le \bar E_s,\quad E_{s,T-1}\ge E^{fin}_s,
\label{eq:soc}\\
\underline P^{ch}_s \xi^+_{s,t}&\le P^{ch}_{s,t}\le \bar P^{ch}_s \xi^+_{s,t},&
\underline P^{dis}_s \xi^-_{s,t}&\le P^{dis}_{s,t}\le \bar P^{dis}_s \xi^-_{s,t},&
\xi^+_{s,t}+\xi^-_{s,t}&\le 1,
\label{eq:sto-bin}
\end{align}
以及日循环上限
$\sum_t\big(\tfrac{\Delta t}{\eta^{dis}_s}P^{dis}_{s,t}+\eta^{ch}_s\Delta t
P^{ch}_{s,t}\big)\le 2 N^{cycle}_s \bar E_s$（吞吐能量口径，充放两侧各计一次，
故右端为两倍循环乘容量）。分效率缺省时按 $\eta^{ch}=\eta^{dis}=\sqrt\eta$
均分往返效率。

\begin{remark}[无指示变量时的松弛]
当 $\xi^\pm$ 不引入模型时，充放互斥只能靠目标中的效率损耗与报价结构"软"达成；
若放电报价低于放电成本折算，LP 松弛可能出现同时充放的虚假套利。实现以
$P^{ch}_{s,t}+P^{dis}_{s,t}\le\max(\bar P^{ch}_s,\bar P^{dis}_s)$ 作软互斥
（见第 4 章声明），严格互斥需开 \fld{use\_binary\_indicators}。
\end{remark}

\subsubsection{直流线路}

\begin{equation}
\underline P^{dc}_d\le P^{dc}_{d,t}\le \bar P^{dc}_d,\qquad
P^{dc}_{d,t}-P^{dc}_{d,t-1}\le R^{dc,+}_d,\quad
P^{dc}_{d,t-1}-P^{dc}_{d,t}\le R^{dc,-}_d,
\label{eq:dcline}
\end{equation}
直流功率可正可负（双向送电），在功率平衡中作为注入项、在网络约束中经
PTDF 差 $\mathrm{PTDF}_{l,to}-\mathrm{PTDF}_{l,from}$ 计入。

\subsection{分段线性化}
\label{sec:pwl}

\begin{definition}[分段线性能量成本]
机组能量成本取报价曲线的下凸分段线性函数
$C_g(P)=C_g(\underline P_g)+\sum_k c_{g,k} p_k$（$P=\underline P_g+\sum_k p_k$，
$0\le p_k\le q_k$），段价单调不减 $c_{g,1}\le c_{g,2}\le\cdots$ 时
最小化问题自动按段序填满（凸 PWL 无歧义）；段价非单调时，\eqref{eq:gen-coupling}
的段上界耦合 $p_{g,k,t}\le q_{g,k}u_{g,h(t)}$ 仍是有效约束，但段的使用顺序由
优化决定，报价曲线应为凸递增以符合市场规则 \cite{southern2025,wood2013}。
\end{definition}

\begin{remark}
二次成本 $aP^2+bP$ 的常用三段等宽线性化：段宽 $w=(\bar P-\underline P)/3$，
第 $k$ 段平均边际成本 $c_k=b+2a\big(\underline P+(2k-1)w/2\big)$——这正是
算例构造器 \srcpath{src/scuc/case\_builder.cpp:build\_ieee118\_case} 中
\texttt{make\_segments} lambda 采用的公式。
\end{remark}

\subsection{直流潮流与 PTDF}
\label{sec:dcflow}

直流潮流假设：无损、电压幅值恒定 1.0 pu、相角差小，得线性系统
\begin{equation}
B\,\theta = p^{inj},
\end{equation}
其中 $B\in\mathbb R^{n_b\times n_b}$，$B_{ff}=B_{tt}=\sum_{l\ni f,t} 1/x_l$、
$B_{ft}=B_{tf}=-1/x_l$（仅计入在运支路）。删去平衡母线（取 0 号）行列得
约化矩阵 $B_{red}$，可逆时支路潮流为
\begin{equation}
F_l=\frac{\theta_{f_l}-\theta_{t_l}}{x_l}
=\sum_{b}\underbrace{\frac{1}{x_l}\big(B^{-1}_{red}\big)_{f_l b}
-\frac{1}{x_l}\big(B^{-1}_{red}\big)_{t_l b}}_{\displaystyle \mathrm{PTDF}_{l,b}}
\ p^{inj}_b ,
\label{eq:ptdf}
\end{equation}
（涉及平衡母线的行列按 0 处理）。$\mathrm{PTDF}_{l,b}$ 即母线 $b$ 单位注入
在支路 $l$ 上引起的潮流。网络安全约束为
\begin{equation}
-\bar F_l\le \sum_b \mathrm{PTDF}_{l,b}\, p^{inj}_b \le \bar F_l,\quad\forall l,t,
\label{eq:linelim}
\end{equation}
$p^{inj}_b$ 由母线 $b$ 处机组出力、可再生、储能净放、直流注入与负荷代数和
构成。监视断面（section）潮流定义为成员线路潮流的加权和，其灵敏度矩阵即
PTDF 行的同一加权组合：
\begin{equation}
\mathrm{SECPTDF}_{s,b}=\sum_{l} w_{s,l}\,\mathrm{PTDF}_{l,b},\qquad
-\bar F^{rev}_s\le \sum_b \mathrm{SECPTDF}_{s,b}\,p^{inj}_b\le \bar F^{fwd}_s.
\label{eq:section}
\end{equation}

\begin{remark}[误差与适用条件]
直流近似忽略网损与无功，电压/相角越界无约束表达；径向退化（$B_{red}$ 病态
或奇异）时 PTDF 无定义。潮流精度对高压输电网典型误差在百分之几量级
\cite{stott2009}；需要精确潮流或 N-1 校验时应外接交流潮流/安全分析程序。
\end{remark}

\subsection{MILP 松弛与紧化}
\label{sec:tightening}

SCUC 的可解性取决于 LP 松弛强度。标准结论（证明从略，见所引文献）：

\begin{proposition}[启停转移的凸包]
单机组三二进制 $(u,y,z)$ 在转移逻辑 \eqref{eq:logic} 与最小启停
\eqref{eq:minupdown} 下的整数可行域，存在由其紧 formulation 给出的 LP 完全
描述 \cite{rajan2005,morales2013}；附加大尺度聚合割（对称破缺、累计段耦合、
爬坡透视割）可显著收紧多机组松弛 \cite{ostrowski2012}。
\end{proposition}

常用紧化手段（均为 LP-valid，即不改变整数可行域）：
\begin{itemize}
  \item \textbf{爬坡透视割}：把 \eqref{eq:ramp} 中 big-M 系数 $\bar P_g$
        收紧为 $\min(\underline P_g+60\Delta t R^+_g,\bar P_g)$，启动时段
        出力上界从 $\bar P_g$ 降到 $\underline P_g+R^+$ 量级，典型收紧
        2--4 倍；
  \item \textbf{转移壳割}：$y_{g,h}\le u_{g,h}$、$z_{g,h}\le 1-u_{g,h}$、
        $y_{g,h}\le 1-u_{g,h-1}$、$z_{g,h}\le u_{g,h-1}$ 给出转移关系的
        凸包刻画；
  \item \textbf{备用覆盖割}：备用需求与单机备用能力按承诺变量聚合，
        $\sum_g \kappa_g u_{g,h(t)}\ge$ 需求，迫使分数解也携带足够在线能力；
  \item \textbf{对称破缺}：参数相同的机组按编号序施加 $u_{g_{j+1},h}\le
        u_{g_j,h}$，削减同构分支。
\end{itemize}

这些割平面在主问题的根节点一次性加入（预构型），其有效性不依赖求解器内部
割生成；代价是根 LP 行数膨胀，故大规模实例需按规模阈值取舍（实现见第 4 章）。

\subsection{SCED 与节点边际电价（LMP）}

\subsubsection{SCED：固定组合的 LP 再调度}

MILP 的整数解固定后（$u,y,z$ 取整固定），SCUC 退化为线性规划——安全约束
经济调度（SCED）。固定组合消除了 MILP 对偶不存在的问题，使 LP 对偶变量成为
有经济解释的价格信号。三阶段分工（SCUC 定组合 → SCED 定出力 → LMP 定价格）
是§2.6.4--2.6.6 口径的市场出清标准流程 \cite{southern2025}。

\subsubsection{LMP 的分解}

\begin{definition}[LMP]
节点边际电价是功率平衡约束与网络约束对偶变量的组合：设 $\lambda_t$ 为时段
$t$ 功率平衡 \eqref{eq:balance} 的对偶（系统边际价格，SMP），
$\mu^+_{l,t},\mu^-_{l,t}\ge 0$ 为线路正/反向限额 \eqref{eq:linelim} 的对偶，
则母线 $b$ 的节点电价为
\begin{equation}
\mathrm{LMP}_{b,t}=\lambda_t+\sum_{l}\big(\mu^-_{l,t}-\mu^+_{l,t}\big)
\,\mathrm{PTDF}_{l,b},
\label{eq:lmp}
\end{equation}
即\textbf{能量分量}（$\lambda_t$，全网相同）与\textbf{阻塞分量}之和
\cite{schweppe1988,hogan1992}。无损直流模型下不含网损分量。
\end{definition}

\begin{remark}[$\delta$ 邻域定价]
MILP 最优组合下的 SCED 对偶尔因退化基或分段拐点处次梯度不唯一而不稳定；
§2.6.5.5 的邻域重调度口径把在线机组出力限制在 SCED 值的
$[(1-\delta),(1+\delta)]$ 倍（并夹在 $[\underline P_g,\bar P_g]$）内重解 LP，
以该 LP 的对偶作为发布价格 \cite{southern2025}。这属于定价规则层面的工程
约定，而非新的优化理论。
\end{remark}

\subsubsection{与场景/安全校验的关系}

\begin{gapnote}
规则意义上的完整 SCUC 还包含：（i）N-1 预想事故集下的网络安全校验（事故态
PTDF/LODF 约束或迭代式安全校核）；（ii）多场景/随机 UC 以覆盖可再生与负荷
预测误差；（iii）交流潮流校验与无功/电压约束。当前代码\textbf{均未实现}：
网络约束仅基态直流潮流，输入为单条确定性预测曲线（\fld{profiles}），无场景
维度；\fld{startup\_cost\_warm}/\fld{startup\_cost\_cold} 与启停过程
（\fld{ud\_periods}/\fld{dd\_periods}）等规则要素已解析未建模。这些差异在
第 4 章转写层逐条声明。
\end{gapnote}

\subsection{与实现的对应}

\begin{footnotesize}
\begin{longtable}{@{}L{5.6cm}L{8.6cm}@{}}
\toprule
\textbf{理论条目} & \textbf{实现状态}\\
\midrule
\endhead
目标函数 \eqref{eq:scuc-obj}（能量+空载+启动+备用+罚） &
\implfull{src/scuc/scuc.cpp:build\_formulation}（启动成本只用
\texttt{startup\_cost}；温/冷启动字段未消费）\\
功率平衡 \eqref{eq:balance} &
\implfull{src/scuc/scuc.cpp:build\_formulation}（含切负荷/过剩松弛变量）\\
备用需求与耦合 \eqref{eq:reserve}\eqref{eq:reserve-coupling} &
\implfull{src/scuc/scuc.cpp:build\_formulation}（含负备用与 PFR 容量形式）\\
出力耦合与分段 \eqref{eq:gen-coupling} &
\implfull{src/scuc/scuc.cpp:build\_formulation}\\
启停逻辑 \eqref{eq:logic}、最小启停 \eqref{eq:minupdown} &
\implfull{src/scuc/scuc.cpp:build\_formulation}（后视形式；时界起点截断，
初始状态历史不进约束）\\
爬坡 \eqref{eq:ramp} &
\implfull{src/scuc/scuc.cpp:build\_formulation}（big-M 取 $\bar P_g$；
紧系数形式作为 Family R 割平面另行加入）\\
储能 \eqref{eq:soc}\eqref{eq:sto-bin} 与循环上限 &
\implfull{src/scuc/scuc.cpp:build\_formulation}（$\xi^\pm$ 由
\fld{use\_binary\_indicators} 控制；否则软互斥）\\
直流线路 \eqref{eq:dcline} &
\implfull{src/scuc/scuc.cpp:build\_formulation}\\
PTDF \eqref{eq:ptdf} 与线路限额 \eqref{eq:linelim} &
\implfull{src/scuc/scuc.cpp:compute\_ptdf}（平衡母线固定为 0，
Eigen FullPivLU 求逆；限额 $>10^5$ MW 的支路不加约束）\\
断面约束 \eqref{eq:section} &
\implfull{src/scuc/scuc.cpp:build\_formulation}（SEC\_PTDF 为 PTDF 行加权和）\\
预构型割平面族（\ref{sec:tightening} 节） &
\implpart{实现 6/A/T/H/R/G/P/Q/C/V/11 共 11 族；$n_g\times T_c>800$ 时
关闭 T/H/R/P/Q 五族，见 \srcpath{src/scuc/scuc.cpp:build\_formulation}}\\
SCED 固定组合 LP &
\implfull{src/scuc/scuc.cpp:solve\_sced\_stage}\\
LMP 分解 \eqref{eq:lmp} 与 $\delta$ 邻域 &
\implfull{src/scuc/scuc.cpp:solve\_lmp\_stage}（对偶取自
\srcpath{src/scuc/scuc.cpp:run\_lp\_with\_duals}）\\
热/温/冷态启动成本 & \implpart{字段已解析（\srcpath{src/scuc/scuc.cpp:scuc\_from\_json}），
目标中一律按热态成本 \texttt{startup\_cost} 计费}\\
启停过程出力曲线（§2.6.3.8 UD/DD） & \implnone（\fld{ud\_periods}/\fld{dd\_periods} 已解析未消费）\\
N-1 安全校验、多场景/随机 UC、交流潮流 & \implnone\\
\bottomrule
\end{longtable}
\end{footnotesize}

\subsection{参考文献}

\begin{thebibliography}{9}\small
\bibitem{southern2025} 南方区域电力市场管理委员会，《南方区域电力市场现货交易
规则（2025 V1.0）》，§2.6 日前市场出清模型（\texttt{src/scuc/scuc.cpp} 与
\texttt{include/mipsolvers/scuc/scuc.hpp} 头部注释所引规则文本）。
\bibitem{wood2013} A.~J. Wood, B.~F. Wollenberg, G.~B. Shebl\'e,
\emph{Power Generation, Operation, and Control}, 3rd ed., Wiley, 2013.
\bibitem{schweppe1988} F.~C. Schweppe, M.~C. Caramanis, R.~D. Tabors,
R.~E. Bohn, \emph{Spot Pricing of Electricity}, Kluwer Academic, 1988.
\bibitem{hogan1992} W.~W. Hogan, ``Contract networks for electric power
transmission,'' \emph{Journal of Regulatory Economics} 4(3):211--242, 1992.
\bibitem{rajan2005} D.~Rajan, S.~Takriti, ``Minimum up/down polytopes of the
unit commitment problem with start-up costs,'' IBM Research Report RC23628,
2005.
\bibitem{ostrowski2012} J.~Ostrowski, M.~F. Anjos, A.~Vannelli, ``Tight mixed
integer linear programming formulations for the unit commitment problem,''
\emph{IEEE Transactions on Power Systems} 27(1):39--46, 2012.
\bibitem{morales2013} G.~Morales-Espa\~na, J.~M. Latorre, A.~Ramos, ``Tight
and compact MILP formulation for the thermal unit commitment problem,''
\emph{IEEE Transactions on Power Systems} 28(4):4897--4908, 2013.
\bibitem{stott2009} B.~Stott, J.~Jardim, O.~Alsa\c{c}, ``DC power flow
revisited,'' \emph{IEEE Transactions on Power Systems} 24(3):1290--1300, 2009.
\end{thebibliography}


% source_equivalent_scuc.tex —— scuc 实现转写：JSON → MILP → 三阶段求解
% 本文件为实现转写章节，遵循 docs/modules/README.md 数学模型规范：
% 公式只转写当前源码实际执行的内容；每组公式给出 文件:函数名 锚点；禁止行号。

\section{实现转写：JSON 到 MILP 的完整装配}

\subsection{数据流总览}

模块数据流（括号内为实现函数）：

\begin{enumerate}
  \item JSON 文本 $\to$ \texttt{SCUCInput}
        （\srcpath{src/scuc/scuc.cpp:scuc\_from\_json}）；
  \item \texttt{SCUCInput} $\to$ 变量布局 \texttt{VarIndex}
        （\srcpath{src/scuc/scuc.cpp:setup\_var\_index}）与 PTDF
        （\srcpath{src/scuc/scuc.cpp:compute\_ptdf}）；
  \item 装配 \texttt{engine::MIPModel}
        （\srcpath{src/scuc/scuc.cpp:build\_formulation}），同时生成原始修复
        种子（\srcpath{src/scuc/scuc.cpp:build\_scuc\_primal\_repair\_seed}）、
        UC 提示与分支优先级；
  \item MILP 求解（\srcpath{src/scuc/scuc.cpp:run\_milp}）与结果提取
        （\srcpath{src/scuc/scuc.cpp:extract\_result}），得 SCUC 阶段结果；
  \item 可选 SCED（\srcpath{src/scuc/scuc.cpp:solve\_sced\_stage}）与 LMP
        （\srcpath{src/scuc/scuc.cpp:solve\_lmp\_stage}）；
  \item \texttt{SCUCOutput} $\to$ JSON
        （\srcpath{src/scuc/scuc.cpp:scuc\_output\_to\_json}）。
\end{enumerate}

以下转写中，"$\le$ 行"对应 \texttt{build\_formulation} 内 lambda
\texttt{add\_le}（写入 \texttt{lp.A}/\texttt{lp.b}），"$=$ 行"对应
\texttt{add\_eq}（写入 \texttt{lp.Aeq}/\texttt{lp.beq}）；两个 lambda 均过滤
非有限或零系数项，右端非有限时置 0
（\srcpath{src/scuc/scuc.cpp:build\_formulation}）。全局数值零阈值
\texttt{kEps}$=10^{-9}$。

\subsection{JSON 解析}
\label{sec:impl-json}

\srcpath{src/scuc/scuc.cpp:scuc\_from\_json} 用 nlohmann::json 解析。逐字段
\texttt{contains} 判在，缺失字段保留结构体默认值（第 6 章"默认值"列）；
数组类顶层键（\fld{generators}、\fld{branches} 等）缺失即为空数组。
\fld{num\_buses} 缺省时取所有元件 \texttt{bus}/\texttt{from}/\texttt{to}
最大值加 1（下限 1）。异常口径：JSON 语法错误抛
\texttt{nlohmann::json::parse\_error}，字段类型不匹配抛
\texttt{nlohmann::json::type\_error}，解析层\textbf{不做}物理合法性校验
（如 $\underline P\le\bar P$、母线索引越界等在装配期静默跳过或截断，见各节）。

\subsection{变量布局与索引空间}
\label{sec:impl-varindex}

\texttt{VarIndex} 把 MILP 变量排成单一平坦向量，块内按
$\mathrm{col}=\mathrm{base}+u\times T+t$（启停类用 $T_c$）编址。装配顺序与
块大小（\srcpath{src/scuc/scuc.cpp:setup\_var\_index}）：

\begin{center}
\begin{footnotesize}
\begin{longtable}{@{}llll@{}}
\toprule
\textbf{块} & \textbf{符号} & \textbf{大小} & \textbf{类型/说明}\\
\midrule
\endhead
\texttt{SU}/\texttt{SD}/\texttt{IG} & $y,z,u$ & $3\times n_g T_c$ & 二进制 $\{0,1\}$\\
\texttt{SEG}$_k$ & $p_{g,k,t}$ & $K\times n_g T$ & 连续，$\ge 0$；上界逐段设定\\
\texttt{PG} & $P_{g,t}$ & $n_g T$ & 连续，$\ge 0$\\
\texttt{RG\_spin}/\texttt{RG\_reg\_up}/\texttt{RG\_reg\_down} & $R$ & $3\times n_g T$ & 连续，$\ge 0$\\
\texttt{PF} & — & $n_l T$ & \textbf{已分配未使用}：无任何约束或目标引用该块\\
\texttt{PW}/\texttt{PPV} & $P_{w,t},P_{s,t}$ & $n_w T+n_{pv}T$ & 连续；上界见可再生节\\
\texttt{WIND\_CURT}/\texttt{SOL\_CURT} & $P^{wc},P^{sc}$ & $n_w T+n_{pv}T$ & 仅当 $M_2>10^{-9}$ 分配，否则块长 0\\
\texttt{PSTO\_IN}/\texttt{PSTO\_OUT}/\texttt{SOC} & $P^{ch},P^{dis},E$ & $3\times n_s T$ & 连续，$\ge 0$\\
\texttt{XI\_CH}/\texttt{XI\_DIS} & $\xi^+,\xi^-$ & $2\times n_s T$ & 二进制；仅当任一储能 \fld{use\_binary\_indicators}=true 分配\\
\texttt{PDC} & $P^{dc}_{d,t}$ & $n_{dc}T$ & 连续，上下界逐线设定（可为负）\\
\texttt{LOAD\_SHED}/\texttt{GEN\_CURT} & $P^{shed}_t,P^{curt}_t$ & $2T$ & 连续，$\ge 0$（\texttt{n\_areas=1} 硬编码）\\
\texttt{SL\_LINE\_POS}/\texttt{NEG} & $\sigma^+,\sigma^-$ & $2\times n_l T$ & 线路松弛，上界 $10^6$\\
\texttt{SL\_SEC\_POS}/\texttt{NEG} & — & $2\times n_{sec}T$ & 断面松弛，上界 $10^6$\\
\bottomrule
\end{longtable}
\end{footnotesize}
\end{center}

所有连续变量默认 $\mathrm{lb}=0$、$\mathrm{ub}=10^{20}$，再按约束族覆盖上界；
二进制块同时登记到 \texttt{mip.binary\_idx}
（\srcpath{src/scuc/scuc.cpp:build\_formulation}）。时间结构：
$\mathrm{iph}=\max(1,\mathrm{round}(1/\Delta t))$，$T_c=\max(1,T/\mathrm{iph})$，
耦合映射 $h(t)=\lfloor t/\mathrm{iph}\rfloor$。

\subsection{PTDF 计算}

对应源码为 \srcpath{src/scuc/scuc.cpp:compute\_ptdf}。步骤：
\begin{enumerate}
  \item 对在运支路（\texttt{in\_service} 且端点在 $[0,n_b)$）累加
        $B$ 矩阵：$b_l=1/x_l$（$|x_l|\le 10^{-6}$ 时按 $x_l=10^{-3}$ 兜底），
        $B_{ff}\mathrel{+}=b_l$，$B_{tt}\mathrel{+}=b_l$，$B_{ft},B_{tf}\mathrel{-}=b_l$；
  \item 平衡母线硬编码为 0 号；删其行列得 $B_{red}$，用
        \texttt{Eigen::FullPivLU} 判可逆并求逆；不可逆时 $B^{-1}_{red}$ 取零矩阵
        （\textbf{回退}：等价于全部 PTDF 为 0，网络约束退化为
        $-\sigma^-\le \mathrm{load\_ptdf}$ 形式的无约束行，见下节）；
  \item $\mathrm{PTDF}_{l,n}=b_l\big(\tilde B^{-1}_{f_l n}-\tilde B^{-1}_{t_l n}\big)$，
        涉及平衡母线的项按 0 计。
\end{enumerate}
$n_l=0$ 或 $n_b\le 1$ 时返回空矩阵，网络约束整块跳过。PTDF 为稠密
$n_l\times n_b$ 矩阵。

\subsection{目标函数装配}

对应源码为 \srcpath{src/scuc/scuc.cpp:build\_formulation}。目标向量
\texttt{lp.c} 各分块系数（$\Delta t$ = \fld{period\_length\_hr}，
$\mathrm{iph}$ 如上）：

\begin{align}
\texttt{SEG}_k(g,t)&:\ c_{g,k}\,\Delta t &&\text{（缺段取 }0\text{）}\nonumber\\
\texttt{PG}(g,t)&:\ \mathrel{+}= c^{wf}\Delta t &&\text{仅当 }c^{wf}>10^{-9}\nonumber\\
\texttt{RG}_{\cdot}(g,t)&:\ c^{\cdot}_g\,\Delta t &&\text{三类备用报价}\nonumber\\
\texttt{SU}(g,h)&:\ \mathrel{+}= C^U_g &&\text{恒用 \fld{startup\_cost}（热态值）}\label{eq:impl-su-cost}\\
\texttt{IG}(g,h)&:\ \mathrel{+}= C^{NL}_g/\mathrm{iph} &&\label{eq:impl-nl-cost}\\
\texttt{LOAD\_SHED}(t)&:\ c^{shed}\Delta t,\ \mathrm{ub}=10^6 &\quad
\texttt{GEN\_CURT}(t)&:\ c^{curt}\Delta t,\ \mathrm{ub}=10^6\nonumber\\
\texttt{SL}_{\cdot}&:\ M_1,\ \mathrm{ub}=10^6 &&\text{线路与断面四块同}\nonumber\\
\texttt{WIND\_CURT},\texttt{SOL\_CURT}&:\ M_2\Delta t,\ \mathrm{ub}=10^6 &&\text{仅当 }M_2>10^{-9}\nonumber\\
\texttt{PSTO\_OUT}(s,t)&:\ \mathrel{+}=\lambda^{dis}_s\Delta t;\quad
\texttt{PSTO\_IN}(s,t)\ \mathrel{+}=\lambda^{ch}_s\Delta t &&\text{仅当 }|\lambda|>10^{-9}\nonumber
\end{align}

\textbf{口径声明}：
\begin{itemize}
  \item \eqref{eq:impl-su-cost}：\fld{startup\_cost\_warm}/\fld{startup\_cost\_cold}
        与 \fld{hot\_start\_threshold\_hr}/\fld{warm\_start\_threshold\_hr}
        已解析但\textbf{未进入}目标或约束；所有启动均按 \fld{startup\_cost} 计费；
  \item \eqref{eq:impl-nl-cost}：空载成本系数为 $C^{NL}_g/\mathrm{iph}$。因
        $\mathrm{iph}\approx 1/\Delta t$，这等价于"每组合时段收
        $C^{NL}_g\cdot\Delta t$"；当 $\Delta t=1$ 时与 \$/h 口径一致，
        $\Delta t<1$ 时全时界空载成本合计为 $C^{NL}_g\Delta t\, T_c$，
        即按时界小时数的 $\Delta t$ 倍计收（如实转写，非 \$/h 严格积分）；
  \item 输电费 $c^{wf}$ 加在\textbf{全部}机组出力上，不按省间/断面区分。
\end{itemize}

\subsection{约束族：机组出力与分段}
\label{sec:impl-gen}

对应源码为 \srcpath{src/scuc/scuc.cpp:build\_formulation}。对每台机组 $g$、
每调度时段 $t$（$h=h(t)$）：
\begin{align}
P_{g,t}-\underline P_g u_{g,h}-\textstyle\sum_k p_{g,k,t}&=0
  &&\text{（}\underline P_g\le 10^{-9}\text{ 时省略 }u\text{ 项）}\label{eq:impl-pgdef}\\
P_{g,t}-\bar P_g u_{g,h}&\le 0 \label{eq:impl-pgub}\\
-P_{g,t}+\underline P_g u_{g,h}&\le 0
  &&\text{（仅当 }\underline P_g>10^{-9}\text{）}\label{eq:impl-pglb}\\
p_{g,k,t}&\le q_{g,k} &&\text{（变量上界；缺段 }q=0\text{）}\label{eq:impl-segub}\\
p_{g,k,t}-q_{g,k}u_{g,h}&\le 0
  &&\text{（仅当 }q_{g,k}>10^{-9}\text{）}\label{eq:impl-segcoupling}
\end{align}
$q_{g,k}$ 装配时取 $\max(0,\cdot)$ 截断。注意 \eqref{eq:impl-pgdef} 与
\eqref{eq:impl-segub} 联立使离网机组（$u=0$）出力恒为
$P=0+\sum p_k$，而 \eqref{eq:impl-segcoupling} 把各段压到 0。

\subsection{约束族：启停逻辑、最小启停与次数上限}

对应源码为 \srcpath{src/scuc/scuc.cpp:build\_formulation}。对每台机组 $g$、
每组合时段 $h$：
\begin{align}
u_{g,0}-y_{g,0}+z_{g,0}&=u_{g,0^-}
  &&\text{（}u_{g,0^-}=\mathrm{clamp}(\fld{initial\_status.commitment},0,1)\text{，缺省 }0\text{）}\label{eq:impl-trans0}\\
u_{g,h}-u_{g,h-1}-y_{g,h}+z_{g,h}&=0,\quad h\ge 1 \label{eq:impl-trans}\\
y_{g,h}+z_{g,h}&\le 1 \label{eq:impl-excl}\\
\sum_{\tau=\max(0,h-TU_g+1)}^{h} y_{g,\tau}-u_{g,h}&\le 0
  &&\text{（仅当 }TU_g=\lceil T^U_g\rceil>0\text{）}\label{eq:impl-minup}\\
\sum_{\tau=\max(0,h-TD_g+1)}^{h} z_{g,\tau}+u_{g,h}&\le 1
  &&\text{（仅当 }TD_g>0\text{）}\label{eq:impl-mindn}\\
\textstyle\sum_h y_{g,h}&\le \fld{max\_startups}
  &&\text{（仅当 }>0\text{）}\label{eq:impl-maxsu}\\
\textstyle\sum_h z_{g,h}&\le \fld{max\_shutdowns}
  &&\text{（仅当 }>0\text{）}\label{eq:impl-maxsd}
\end{align}
\fld{must\_run}=true 的机组固定 $u_{g,h}\equiv 1$（变量上下界同取 1）。
\textbf{边界声明}：\eqref{eq:impl-minup}/\eqref{eq:impl-mindn} 的求和在时界
起点截断于 $\tau=0$，初始在/离网已持续时间（\fld{time\_in\_state}）
\textbf{不进入} MILP 约束，仅被原始修复启发式消费（见
\ref{sec:impl-repair} 节）。

\subsection{约束族：爬坡}

对应源码为 \srcpath{src/scuc/scuc.cpp:build\_formulation}。记
$R^+=\fld{ramp\_up\_mw\_min}\cdot 60\Delta t$、$R^-=\fld{ramp\_dn\_mw\_min}\cdot
60\Delta t$（MW/时段）：
\begin{align}
P_{g,0}-\bar P_g y_{g,0}&\le P_{g,0^-}+R^+, &
-P_{g,0}-\bar P_g z_{g,0}&\le -P_{g,0^-}+R^-, \label{eq:impl-ramp0}\\
P_{g,t}-P_{g,t-1}-\bar P_g y_{g,h}&\le R^+, &
P_{g,t-1}-P_{g,t}-\bar P_g z_{g,h}&\le R^-,\quad t\ge 1, \label{eq:impl-ramp}
\end{align}
其中 $P_{g,0^-}=\max(0,\fld{initial\_status.dispatch})$（缺省 0）。启动/停机
项的 big-M 系数为 $\bar P_g$；更紧的 $\min(\underline P_g+R^\pm,\bar P_g)$
系数形式以割平面族 R 另行加入（\ref{sec:impl-cuts} 节），不替换本族。

\subsection{约束族：可再生出力}

对应源码为 \srcpath{src/scuc/scuc.cpp:build\_formulation}；预测值由
\srcpath{src/scuc/scuc.cpp:wind\_at}/\srcpath{src/scuc/scuc.cpp:solar\_at}
给出（曲线为绝对 MW；缺行/缺列回退 \fld{pmax}）。记
$F_{w,t}=\fld{profiles.wind}[w][t]$，$\alpha_n=$ \fld{renewable\_min\_output\_coeff}：
\begin{itemize}
  \item $M_2>10^{-9}$（弃电变量启用）：
        \begin{equation}
        P_{w,t}+P^{wc}_{w,t}=\max(0,F_{w,t}),\quad
        P_{w,t}\ge\max(0,\alpha_n F_{w,t}),\quad
        0\le P^{wc}_{w,t}\le\max(0,F_{w,t}),
        \label{eq:impl-renew-curt}
        \end{equation}
        此时 $P_{w,t}$ 上界放开为 $10^{20}$，由等式约束封闭；光伏同构；
  \item $M_2\le 10^{-9}$：无弃电变量，直接夹界
        $\max(0,\alpha_n F_{w,t})\le P_{w,t}\le\max(0,F_{w,t})$；输出层的弃电量
        由 $\max(0,F-P)$ 事后重算（\srcpath{src/scuc/scuc.cpp:extract\_result}）。
\end{itemize}

\subsection{约束族：储能}

对应源码为 \srcpath{src/scuc/scuc.cpp:build\_formulation}。每台储能 $s$ 的
装配期常量：
\begin{align*}
\eta^{ch}&=\begin{cases}\mathrm{clamp}(\fld{eta\_charge},0.01,1),&\fld{eta\_charge}>0\\
\sqrt{\mathrm{clamp}(\fld{efficiency},0.1,1)},&\text{否则}\end{cases}
&&\eta^{dis}\text{ 同构}\\
\underline E&=\big(\fld{soc\_min}\ge 0\;?\;\mathrm{clamp}(\fld{soc\_min},0,1):0.10\big)\cdot\bar E
&&\bar E=\max(0,\fld{energy\_capacity\_mwh})\\
E_0&=\mathrm{clamp}\big(\fld{initial\_status.storage\_soc}\ (\le 1+10^{-8}\text{ 按 pu，否则按 MWh})\ \text{或}\ \fld{soc\_init}\cdot\bar E,\ 0,\bar E\big)\\
E^{fin}&=\big(\fld{soc\_final}\ge 0\;?\;\mathrm{clamp}(\fld{soc\_final},0,1)\cdot\bar E:E_0\big)
\end{align*}
每时段约束：
\begin{align}
\underline E\le E_{s,t}&\le \bar E &&\text{（变量界）}\label{eq:impl-socbnd}\\
E_{s,0}-\eta^{ch}\Delta t\,P^{ch}_{s,0}+\frac{\Delta t}{\eta^{dis}}P^{dis}_{s,0}&=E_0 \label{eq:impl-soc0}\\
E_{s,t}-E_{s,t-1}-\eta^{ch}\Delta t\,P^{ch}_{s,t}+\frac{\Delta t}{\eta^{dis}}P^{dis}_{s,t}&=0,\quad t\ge 1 \label{eq:impl-socdyn}\\
-E_{s,T-1}&\le -E^{fin} \label{eq:impl-socfin}\\
\sum_{t}\Big(\frac{\Delta t}{\eta^{dis}}P^{dis}_{s,t}+\eta^{ch}\Delta t P^{ch}_{s,t}\Big)&\le 2N^{cycle}_s\bar E
&&\text{（仅当 }\fld{cycle\_limit}>0\text{ 且 }\bar E>0\text{）}\label{eq:impl-cycle}
\end{align}
充放互斥两条路径：
\begin{itemize}
  \item 二进制模式（本机 \fld{use\_binary\_indicators}=true 且全局块已分配）：
        \begin{align}
        -P^{ch}_{s,t}+\underline P^{ch}\xi^+_{s,t}&\le 0\ (\text{仅当 }\underline P^{ch}>0),&
        P^{ch}_{s,t}-\bar P^{ch}\xi^+_{s,t}&\le 0,\nonumber\\
        -P^{dis}_{s,t}+\underline P^{dis}\xi^-_{s,t}&\le 0\ (\text{仅当 }\underline P^{dis}>0),&
        P^{dis}_{s,t}-\bar P^{dis}\xi^-_{s,t}&\le 0,\nonumber\\
        \xi^+_{s,t}+\xi^-_{s,t}&\le 1; &&\label{eq:impl-xi}
        \end{align}
  \item 否则（\textbf{LP 松弛口径}）：仅变量上界
        $P^{ch}\le\bar P^{ch}$、$P^{dis}\le\bar P^{dis}$ 加软互斥行
        \begin{equation}
        P^{ch}_{s,t}+P^{dis}_{s,t}\le\max(\bar P^{ch},\bar P^{dis}),
        \label{eq:impl-softmutex}
        \end{equation}
        该口径不保证严格互斥（同时充放在 LP 意义下可行），已在第 3 章注记
        声明；部分机组开二进制、部分不开时两类约束按机组各自生效。
\end{itemize}

\subsection{约束族：直流线路}

对应源码为 \srcpath{src/scuc/scuc.cpp:build\_formulation}：
\begin{align}
\fld{pmin}\le P^{dc}_{d,t}&\le\fld{pmax} &&\text{（变量界）}\nonumber\\
P^{dc}_{d,0}\le\fld{ramp\_up},\quad -P^{dc}_{d,0}&\le\fld{ramp\_dn}
&&\text{（首时段相对 }0\text{）}\label{eq:impl-dc0}\\
P^{dc}_{d,t}-P^{dc}_{d,t-1}\le\fld{ramp\_up},\quad
P^{dc}_{d,t-1}-P^{dc}_{d,t}&\le\fld{ramp\_dn},\quad t\ge 1. \label{eq:impl-dcramp}
\end{align}

\subsection{约束族：系统功率平衡}

对应源码为 \srcpath{src/scuc/scuc.cpp:build\_formulation}，每时段一条等式
（起始行号记入 \texttt{power\_balance\_eq\_start}，供 LMP 阶段定位对偶）：
\begin{equation}
\sum_g P_{g,t}+\sum_w P_{w,t}+\sum_s P_{s,t}
+\sum_{s'}\big(P^{dis}_{s',t}-P^{ch}_{s',t}\big)+\sum_d P^{dc}_{d,t}
+P^{shed}_t-P^{curt}_t=\sum_{d'}\fld{load\_at}(d',t),
\label{eq:impl-balance}
\end{equation}
其中 \srcpath{src/scuc/scuc.cpp:load\_at} 实现"基准 MW × 曲线 pu"，曲线
缺行/缺列回退基准值。

\subsection{约束族：线路潮流（PTDF）}

对应源码为 \srcpath{src/scuc/scuc.cpp:build\_formulation}。\textbf{跳过条件}
（硬编码）：支路 \fld{in\_service}=false，或限额
$\bar F_l=\max(0,\fld{rating\_mw})>10^5$ MW（视为不阻塞，不生成约束行）。
对其余支路、每时段 $t$：
\begin{align}
\sum_b \mathrm{PTDF}_{l,b}\,\mathcal P_{b,t}-\sigma^+_{l,t}&\le
\bar F_l+\Lambda_{l,t} \label{eq:impl-line-fwd}\\
-\sum_b \mathrm{PTDF}_{l,b}\,\mathcal P_{b,t}-\sigma^-_{l,t}&\le
\bar F_l-\Lambda_{l,t} \label{eq:impl-line-rev}
\end{align}
其中 $\mathcal P_{b,t}$ 为母线 $b$ 处可调注入（机组 $P_{g,t}$、风电 $P_{w,t}$、
光伏 $P_{s,t}$、储能 $P^{dis}-P^{ch}$；$|\mathrm{PTDF}|\le 10^{-12}$ 的母线
跳过），直流线以系数 $\mathrm{PTDF}_{l,to}-\mathrm{PTDF}_{l,from}$ 计入，
负荷项移入右端
$\Lambda_{l,t}=\sum_b \mathrm{PTDF}_{l,b}\sum_{d'\in b}\fld{load\_at}(d',t)$。
松弛变量 $\sigma^\pm\ge 0$、上界 $10^6$、目标罚系数 $M_1$。每条行同时被
\texttt{record\_network\_flow\_row}（\srcpath{src/scuc/scuc.cpp:build\_formulation}
内 lambda）登记到 UC 提示的 network\_flow\_* 数组（不含松弛列），供原生 B\&C
的网络承诺割使用。

\subsection{约束族：监视断面}

对应源码为 \srcpath{src/scuc/scuc.cpp:build\_formulation}。断面灵敏度
\begin{equation}
\mathrm{SECPTDF}_{s,:}=\sum_{(l,w_l)\in\fld{line\_weights}} w_l\,\mathrm{PTDF}_{l,:}
\quad(\text{越界线路索引跳过}),
\label{eq:impl-secptdf}
\end{equation}
约束结构同 \eqref{eq:impl-line-fwd}/\eqref{eq:impl-line-rev}，限额分别取
\fld{rating\_fwd\_mw}/\fld{rating\_rev\_mw}（取 $\max(0,\cdot)$），松弛
\texttt{SL\_SEC\_POS/NEG}、罚 $M_1$。仅当 $n_{sec}>0$ 且 PTDF 非空时生成。

\subsection{约束族：备用需求与耦合}

对应源码为 \srcpath{src/scuc/scuc.cpp:build\_formulation}。记
$D_t=\sum_{d'}\fld{load\_at}(d',t)$。每时段系统需求行（$>0$ 才生成）：
\begin{align}
-\textstyle\sum_g R^{spin}_{g,t}&\le -r^{spin}D_t, &
-\textstyle\sum_g R^{up}_{g,t}&\le -r^{up}D_t, &
-\textstyle\sum_g R^{dn}_{g,t}&\le -r^{dn}D_t, \label{eq:impl-resreq}
\end{align}
单机耦合（每 $g,t$ 恒生成，$h=h(t)$）：
\begin{align}
P_{g,t}+R^{spin}_{g,t}+R^{up}_{g,t}-\bar P_g u_{g,h}&\le 0 \label{eq:impl-headroom}\\
R^{dn}_{g,t}-P_{g,t}+\underline P_g u_{g,h}&\le 0 \label{eq:impl-regdn}\\
R^{spin}_{g,t}-(\bar P_g-\underline P_g)u_{g,h}&\le 0 \label{eq:impl-spincap}\\
R^{up}_{g,t}-10\,\fld{ramp\_up\_mw\_min}\,u_{g,h}&\le 0
&&\text{（仅当 }10R^+_{min}>10^{-9}\text{）}\label{eq:impl-rupcap}\\
R^{dn}_{g,t}-10\,\fld{ramp\_dn\_mw\_min}\,u_{g,h}&\le 0
&&\text{（仅当 }10R^-_{min}>10^{-9}\text{）}\label{eq:impl-rdncap}
\end{align}
\textbf{口径声明}：\eqref{eq:impl-rupcap}/\eqref{eq:impl-rdncap} 的系数是
$10\times$ MW/min 爬坡率（即 10 分钟可调容量，单位 MW），
\textbf{不乘} $\Delta t$；当 $\Delta t\ne 1/6$ h 时该口径仍按 10 分钟计。

负备用与一次调频（仅当对应需求 $>0$）：
\begin{align}
\sum_g\big(P_{g,t}-\underline P_g u_{g,h}\big)&\le D_t-r^{neg}D_t
&&\text{（§2.6.3.3 口径）}\label{eq:impl-negres}\\
-\sum_g \alpha^{pfr}_g\bar P_g u_{g,h}&\le -R^{pfr}
&&\text{（§2.6.3.4 容量形式）}\label{eq:impl-pfr}
\end{align}

\subsection{约束族：机组群}

对应源码为 \srcpath{src/scuc/scuc.cpp:build\_formulation}。对每个非空群：
\begin{align}
\textstyle\sum_{g\in grp}P_{g,t}&\le \fld{pmax\_t}[t^\ast], &
-\textstyle\sum_{g\in grp}P_{g,t}&\le -\fld{pmin\_t}[t^\ast], \label{eq:impl-grpout}\\
\sum_{g\in grp}\sum_t \Delta t\,P_{g,t}&\le \fld{emax}, &
-\sum_{g\in grp}\sum_t \Delta t\,P_{g,t}&\le -\fld{emin}, \label{eq:impl-grpeng}
\end{align}
其中 $t^\ast=t$（向量长度 $>1$，越界取末元素）或 $0$（长度 1 全时段共用）；
\fld{pmin\_t}/\fld{pmax\_t} 为空则对应侧不生成；\fld{emin}/\fld{emax} 为 0
侧不生成。\fld{generators[].group\_id} 字段与本族无联动（已解析未消费）。

\subsection{预构型市场割平面}
\label{sec:impl-cuts}

\fld{enable\_market\_cuts}=true 时（默认），在装配期一次性加入下列 LP-valid
割平面族（均不改变整数可行域；计数记入结果的 \fld{n\_cuts\_added}）。
对应源码均为 \srcpath{src/scuc/scuc.cpp:build\_formulation}。规模阈值
（硬编码）：$\mathrm{large\_uc}\Leftrightarrow n_g\times T_c>800$，
成立时关闭 T、H、R、P、Q 五族。

\begin{itemize}
  \item \textbf{Family 6（对称破缺）}：按
        $(\mathrm{round}(\bar P),\mathrm{round}(\underline P),\mathrm{round}(T^U),
        \mathrm{round}(T^D))$ 分组，组内按编号序
        $u_{g_{j+1},h}-u_{g_j,h}\le 0$；
  \item \textbf{Family A（累计段耦合）}：对段前缀 $m\ge 1$，
        $\sum_{k\le m}p_{g,k,t}-\big(\sum_{k\le m}q_{g,k}\big)u_{g,h}\le 0$
        （仅当累计量 $>10^{-9}$）；
  \item \textbf{Family T（转移壳，}large\_uc \textbf{时关）}：每 $g,h$ 四行
        $y\le u$、$z\le 1-u$、$y_h\le 1-u_{h-1}$、$z_h\le u_{h-1}$
        （$h=0$ 时以初值 $u_{0^-}$ 代 $u_{-1}$）；
  \item \textbf{Family G（启动团）}：窗口 $W=TU_g+TD_g\ge 2$ 时
        $\sum_{\tau=h}^{h+W-1}y_{g,\tau}\le 1$（$h+W-1<T_c$）；
  \item \textbf{Family H（转移冲突覆盖，}large\_uc \textbf{时关）}：
        $TU_g>1$ 时 $y_{g,s}+z_{g,s'}\le 1$（$s<s'\le s+TU_g-1$）；
        $TD_g>1$ 时对称；
  \item \textbf{Family R（紧爬坡透视，}large\_uc \textbf{时关）}：
        记 $\tilde R^\pm=\min(\bar P_g,60\Delta t R^\pm_g)$、
        $S^+=\min(\bar P_g,\underline P_g+\tilde R^+)$、
        $S^-=\min(\bar P_g,\underline P_g+\tilde R^-)$，对 $t\ge 1$：
        \begin{align}
        P_{g,t}-P_{g,t-1}-\tilde R^+ u_{g,h(t-1)}-S^+ y_{g,h(t)}&\le 0,
        \label{eq:impl-famr}\\
        P_{g,t-1}-P_{g,t}-\tilde R^- u_{g,h(t)}-S^- z_{g,h(t)}&\le 0;
        \nonumber
        \end{align}
  \item \textbf{Family P（启动爬坡多步上界，}large\_uc \textbf{或}
        $T\ne T_c$ \textbf{时关）}：$k=1,\dots,\min(TU_g-1,8)$，
        $\mathrm{excess}=\bar P_g-\underline P_g-k\tilde R^+>10^{-9}$ 时
        \begin{equation}
        P_{g,h}+\mathrm{excess}\cdot y_{g,h-k+1}-\bar P_g u_{g,h}\le 0,
        \quad h\ge k-1;
        \label{eq:impl-famp}
        \end{equation}
  \item \textbf{Family Q（停机爬坡多步上界}，同条件\textbf{）}：
        $P_{g,h}+\mathrm{excess}^-\cdot z_{g,h+k}-\bar P_g u_{g,h}\le 0$
        （$h+k<T_c$）；
  \item \textbf{Family C（系统备用覆盖聚合）}：上方向
        \begin{equation}
        -\sum_g \bar P_g u_{g,h}-\sum_w P_{w,t}-\sum_s P_{s,t}
        -\sum_{s'}(P^{dis}_{s',t}-P^{ch}_{s',t})-\sum_d P^{dc}_{d,t}
        -P^{shed}_t+P^{curt}_t\le -(D_t+r^{up,tot}D_t),
        \label{eq:impl-famc}
        \end{equation}
        其中 $r^{up,tot}=r^{spin}+r^{up}$；下方向对偶地把系数反号、
        以 $\underline P_g$ 代 $\bar P_g$、右端取 $D_t-r^{dn}D_t$；
  \item \textbf{Family V（备用可交付覆盖）}：对四组（需求， 单机能力）
        ——$(r^{spin}+r^{up},\ \bar P-\underline P)$、$(r^{spin},\ \bar P-\underline P)$、
        $(r^{up},\ \min(\bar P-\underline P,10R^+_{min}))$、
        $(r^{dn},\ \min(\bar P-\underline P,10R^-_{min}))$——分别生成
        $-\sum_g \kappa_g u_{g,h}\le -r\,D_t$（能力 $>10^{-9}$ 才入行）；
  \item \textbf{Family 11（储能循环 SOC 界）}：$\bar E>0$ 且右端 $>10^{-9}$ 时
        \begin{equation}
        \sum_t \frac{\Delta t}{\eta^{dis}}P^{dis}_{s,t}\ \le\
        \eta^{ch}\bar P^{ch} T\Delta t+\mathrm{clamp}(\fld{soc\_init},0,1)\bar E-E^{fin}.
        \label{eq:impl-fam11}
        \end{equation}
\end{itemize}

\subsection{UC 提示与分支优先级}

装配末尾填充 \texttt{engine::MIPModel::UCGenHint} 并挂到
\texttt{mip.uc\_hint}（\srcpath{src/scuc/scuc.cpp:build\_formulation}）：
机组数与组合时段数、各二进制/出力块的列索引（$h\times n_g+g$ 排列）、
$\underline P/\bar P$、爬坡（$\max(R^+,R^-)\cdot 60\Delta t$）、四类备用能力、
最小启停取整值、初始状态、逐时段负荷与四类备用需求、分段列与段容量
（仅 $T=T_c$ 时填）、按 $(\underline P,\bar P,R,TU,TD)$ 千分桶聚类的同型机组
分组（含 \texttt{group\_is\_s\_class}：$\underline P=20,\bar P=100,TU=TD=2$
的硬编码识别）、以及 network\_flow\_* 行记录。各
\texttt{certifies\_*\_rows} 标志置真（储能循环行仅当 $n_s>0$）。

分支优先级（\srcpath{src/scuc/scuc.cpp:build\_formulation}）：步长
$S=T_c+1$，$u$ 类优先级 $3S+\rho_g S+(T_c-h)$，$y/z$ 类 $2S+(T_c-h)$；
机组档位 $\rho_g$ 按 $\bar P_g$ 分档（$>600\Rightarrow 6$；$>500\Rightarrow 5$；
$>400\Rightarrow 4$；$>250\Rightarrow 3$；$>150\Rightarrow 2$；
$>100\Rightarrow 1$；否则 0），早时段优先。

\subsection{原始修复种子}
\label{sec:impl-repair}

\fld{enable\_primal\_repair}=true 且环境变量
\texttt{MIPSOLVERS\_DISABLE\_SCUC\_PRIMAL\_REPAIR} 未设置时
（\srcpath{src/scuc/scuc.cpp:scuc\_primal\_repair\_enabled}），
\srcpath{src/scuc/scuc.cpp:build\_scuc\_primal\_repair\_seed} 生成
\texttt{mip.initial\_solution}：

\begin{enumerate}
  \item \textbf{承诺启发式}：按 score 升序排 merit 序，
        $\mathrm{score}=(\fld{must\_run}\,?\,-10^{12}:0)+\mathrm{marginal}
        +(C^{NL}+C^U/\max(1,\lceil T^U\rceil))/\max(1,\bar P)$，
        marginal 取首个正容量段的报价；逐组合时段先满足
        \fld{must\_run} 与锁定期（lock\_on/lock\_off，由
        \fld{time\_in\_state} 折算剩余最小启停期），容量/备用能力不足时按
        merit 序开机（先不越 lock\_off，不足再越），再按逆 merit 序试停
        昂贵机组；启停后重置锁定期；
  \item \textbf{连续恢复}：
        \srcpath{src/scuc/scuc.cpp:recover\_scuc\_continuous\_seed}
        按承诺填充 $u,y,z$、可再生预测、储能回充电量（仅充不放，
        目标为末时段 $\max(\underline E,E^{fin})$）、按 score 序的
        经济出力（含爬坡与备用分配），并调用
        \srcpath{src/scuc/scuc.cpp:scuc\_repair\_network\_slacks}
        两遍扫描把越限行的网络松弛变量顶起；
  \item \textbf{验收与回退}：
        \srcpath{src/scuc/scuc.cpp:scuc\_seed\_satisfies} 以容差 $10^{-5}$
        检查变量界与 $A/Aeq$ 行违反；不可行则回退"全员在线"密种子再验；
        仍不可行时取违反较小者返回（此时种子不保证可行，仅作暖启动参考）。
        设环境变量 \texttt{MIPSOLVERS\_SCUC\_PRIMAL\_REPAIR\_LOG} 可在
        stderr 打印两种种子的最大违反。
\end{enumerate}

\textbf{边界声明}：该种子只是传给 B\&C 后端的初始 incumbent 提示；最终解的
可行性/最优性完全由后端保证，种子不可行不会导致错误结果。

\subsection{三阶段求解流水线}

\begin{itemize}
  \item \textbf{SCUC}：\srcpath{src/scuc/scuc.cpp:solve\_scuc\_stage} 装配后调
        \srcpath{src/scuc/scuc.cpp:run\_milp}。后者注册
        \texttt{StrictHighsBranchAndCutAdapter} 与
        \texttt{NativeBranchAndCutAdapter}（两者均以
        \srcpath{src/scuc/scuc.cpp:make\_scuc\_bc\_options} 构造：映射
        \fld{time\_limit\_sec}/\fld{mip\_gap}/\fld{verbose} 后经
        \srcpath{include/mipsolvers/engine/solver/native/native\_adapters.hpp:make\_strict\_highs\_production\_options}
        补齐生产默认），再注册全部默认适配器；\fld{solver}$\ne$Auto 时设
        \texttt{preferred\_solver}，\fld{allow\_fallback} 控制回退。
  \item \textbf{SCED}：\srcpath{src/scuc/scuc.cpp:solve\_sced\_stage}
        重新装配同一公式，把 $u,y,z$ 四舍五入后固定（lb=ub=取整值，
        类型改连续），清空整数索引，经 \srcpath{src/scuc/scuc.cpp:run\_lp}
        解 LP。
  \item \textbf{LMP}：\srcpath{src/scuc/scuc.cpp:solve\_lmp\_stage} 同样固定
        二进制（参考阶段为 SCED 结果，未运行则为 SCUC 结果），并对在线机组
        出力加 $\delta$ 邻域界
        \begin{equation}
        \max\big(\underline P_g,(1-\delta)P^{sced}_{g,t}\big)\le P_{g,t}\le
        \min\big(\bar P_g,(1+\delta)P^{sced}_{g,t}\big),
        \label{eq:impl-delta}
        \end{equation}
        下界 $>$ 上界 $+10^{-6}$ 时回退 $[\underline P_g,\bar P_g]$；离网机组
        出力固定为 0。LP 经
        \srcpath{src/scuc/scuc.cpp:run\_lp\_with\_duals} 求解——优先
        \texttt{NativeBranchAndCut}（LP 路径）与 \texttt{NativeIPMLP}，因为
        它们返回约束对偶，HiGHS 文件式适配器不返回对偶。
        对偶提取（\srcpath{src/scuc/scuc.cpp:solve\_lmp\_stage}）：
        \begin{equation}
        \lambda_t=\pi\big(n_{ineq}+\texttt{power\_balance\_eq\_start}+t\big),\qquad
        \mathrm{cong}_{b,t}=\sum_l\big(\mu^-_{l,t}-\mu^+_{l,t}\big)\mathrm{PTDF}_{l,b},
        \label{eq:impl-lmp}
        \end{equation}
        其中 $\pi$ 为 \texttt{constraint\_duals} 向量（不等式行在前、等式行
        在后），$\mu^\pm_{l,t}=-\pi(\texttt{line\_\{fwd,rev\}\_row}[lT+t])$；
        $\mathrm{LMP}_{b,t}=\lambda_t+\mathrm{cong}_{b,t}$。
        \textbf{口径声明}：被跳过的线路（退出运行或限额 $>10^5$）行号为 $-1$，
        对阻塞分量贡献 0。
\end{itemize}

\subsection{结果提取与成本分解}

对应源码为 \srcpath{src/scuc/scuc.cpp:extract\_result}。按变量布局回读各矩阵；
线路潮流不由变量回读，而是按 \eqref{eq:ptdf} 以 PTDF $\times$ 节点注入重算
（含负荷与直流项），断面潮流同理由 SEC\_PTDF 重算；弃电量在弃电变量未分配时
按 $\max(0,F-P)$ 重算。成本分解：
\begin{align}
\texttt{energy}&=\textstyle\sum_{g,k,t}c_{g,k}p_{g,k,t}\Delta t, &
\texttt{startup}&=\textstyle\sum_{g,h}C^U_g y_{g,h},\nonumber\\
\texttt{no\_load}&=\textstyle\sum_{g,h}C^{NL}_g u_{g,h}/\mathrm{iph}, &
\texttt{reserve}&=\textstyle\sum_{g,t}(\cdots)\Delta t,\nonumber\\
\texttt{penalty}&=\textstyle\sum_t\big(c^{shed}P^{shed}_t+c^{curt}P^{curt}_t\big)\Delta t,
&&\texttt{total}=\text{前五项之和}.\label{eq:impl-cost}
\end{align}
\textbf{口径声明}：\texttt{total} \textbf{不含}线路/断面松弛罚（$M_1$ 项）、
弃电罚（$M_2$ 项）、储能充放报价与输电费，而求解器返回的
\fld{scuc.objective} 含全部项；存在松弛或这些可选项时
\texttt{objective} $>$ \texttt{cost.total}，二者差异即未分解项之和。
\fld{n\_cuts\_added} 为装配期割平面行数。

\subsection{JSON 序列化}

\begin{itemize}
  \item 输出：\srcpath{src/scuc/scuc.cpp:scuc\_output\_to\_json} 组装
        \texttt{meta}/\texttt{scuc}/\texttt{sced}/\texttt{lmp} 四块
        （\texttt{sced}、\texttt{lmp} 仅在对应开关开启时写出，\textbf{不检查}
        收敛标志）；单阶段块由匿名命名空间
        \srcpath{src/scuc/scuc.cpp:solve\_result\_to\_json} 写出，二维矩阵经
        \srcpath{src/scuc/scuc.cpp:matrix2d\_to\_json} 逐行展开。
        未写入 JSON 的 C++ 字段：\texttt{segment\_dispatch}、
        \texttt{wind\_curtailment}、\texttt{solar\_curtailment}、
        \texttt{section\_flows}；
  \item 输入回写：\srcpath{src/scuc/case\_builder.cpp:scuc\_input\_to\_json}
        把 \texttt{SCUCInput} 全字段导出为可被
        \srcpath{src/scuc/scuc.cpp:scuc\_from\_json} 重新解析的 JSON
        （golden-file 用途）；\fld{time\_in\_state} 仅非空时写出。
\end{itemize}

\subsection{算例构造器（case\_builder）}

对应源码为 \srcpath{src/scuc/case\_builder.cpp}。

\paragraph{时序曲线。}三个静态函数生成曲线（$h_t=(t+0.5)\Delta t$，
$H=T\Delta t$）：
\begin{align}
\fld{make\_load\_profile}&:\ pu_t=f_{min}+(1-f_{min})\max(g_1,g_2),&
g_i&=\exp\Big(-\tfrac12\big(\tfrac{h_t-p_i}{\sigma}\big)^2\Big),\nonumber\\
&&p_1&=\tfrac{10}{24}H,\ p_2=\tfrac{19}{24}H,\ \sigma=\tfrac{2}{24}H,\ f_{min}=0.65;
\label{eq:impl-loadprof}\\
\fld{make\_wind\_profile}&:\ pu_t=\max\big(0.1,\ 1-0.6\,d_t\big),&
d_t&=\exp\Big(-\tfrac12\big(\tfrac{h_t-\frac{13}{24}H}{\frac{4}{24}H}\big)^2\Big);
\label{eq:impl-windprof}\\
\fld{make\_solar\_profile}&:\ pu_t=\begin{cases}
\exp\big(-\tfrac12(\tfrac{h_t-\frac12 H}{\frac18 H})^2\big),&>0.01\\0,&\text{否则}\end{cases}
\label{eq:impl-solarprof}
\end{align}
对应 \srcpath{src/scuc/case\_builder.cpp:make\_load\_profile}、
\srcpath{src/scuc/case\_builder.cpp:make\_wind\_profile}、
\srcpath{src/scuc/case\_builder.cpp:make\_solar\_profile}。负荷曲线为 pu 倍率
（乘 \fld{loads[].p\_mw}），风/光曲线构造后即乘装机容量转为绝对 MW
（构造器侧），并开启 $M_2$（6bus 取 50，其余取 80 \$/MWh）。

\paragraph{算例。}
\begin{itemize}
  \item \srcpath{src/scuc/case\_builder.cpp:build\_3bus\_case}：3 母线 2 机组
        （G1: 30--200 MW，G2: 20--120 MW）、2 支路（限额 150 MW）、单负荷
        180 MW；$K=2$ 段报价；$r^{spin}=0.05$、$r^{up}=0.02$；VOLL=5000、
        VOCC=100、$M_1=10^4$；
  \item \srcpath{src/scuc/case\_builder.cpp:build\_6bus\_case}：Garver 6 母线，
        3 机组（煤 400 MW/气 200 MW/峰 80 MW）、7 支路、4 负荷共 500 MW；
        可选风电（母线 3，100 MW）与储能（母线 5，60 MW/240 MWh，$\eta=0.85$，
        SOC$_{min}$=0.10，循环上限 1，充电报价 $-5$、放电 $+8$ \$/MWh）；
  \item \srcpath{src/scuc/case\_builder.cpp:build\_ieee39\_case}：IEEE 39 母线
        新英格兰系统，10 机组（三段报价、\fld{pfr\_alpha} 0.03--0.06，
        $R^{pfr}=300$ MW）、\textbf{51} 条支路（44 条线路 + G9 升压变
        + 2 条 G10 联络线 + 7 台升压变；注释称 46/51 两种口径，以数组实际
        51 条为准）、17 个非零负荷（26 个条目中 $<1$ MW 的被跳过，
        注释称 21 个），峰值约 6 GW；可选 3 风电（400/300/350 MW）、
        2 光伏（250/200 MW）；
  \item \srcpath{src/scuc/case\_builder.cpp:build\_ieee118\_case}：MATPOWER
        case118 拓扑，54 机组、186 支路（\texttt{static\_assert} 保证）、
        \textbf{98} 个负荷（头文件注释称 99，以数组实际 98 条为准）、
        峰值约 4242 MW；报价由二次成本 $aP^2+bP$ 经
        \texttt{make\_segments} lambda 三段等宽线性化
        （$c_k=b+2a(\underline P+(2k-1)w/2)$，$w=(\bar P-\underline P)/3$）；
        \fld{mip\_gap}=0.003、\fld{time\_limit\_sec}=600、$M_1=10^6$；
        可选 4 风电（500/400/600/350 MW）、3 光伏（300/250/200 MW）。
\end{itemize}

\paragraph{CLI。}\srcpath{src/scuc/case\_builder\_main.cpp:main} 校验
$T\in[1,8760]$、$\Delta t\in(0,24]$，未知算例名报错退出；\fld{--storage}
仅对 6bus 生效（其余算例忽略该开关，无告警）。

\subsection{未消费字段、回退与硬编码阈值清单}

\begin{center}
\begin{footnotesize}
\begin{longtable}{@{}L{4.6cm}L{9.6cm}@{}}
\toprule
\textbf{条目} & \textbf{口径}\\
\midrule
\endhead
\fld{startup\_cost\_warm}/\fld{startup\_cost\_cold}、\fld{hot/warm\_start\_threshold\_hr} &
已解析（\srcpath{src/scuc/scuc.cpp:scuc\_from\_json}），目标与约束均不消费；
启动一律按 \fld{startup\_cost} 计费\\
\fld{ud\_periods}/\fld{dd\_periods} & 已解析，启停过程出力曲线未建模\\
\fld{generators[].group\_id} & 已解析，不参与群约束（群约束只看
\fld{generator\_groups[].gen\_indices}）\\
\fld{initial\_status.time\_in\_state} & 仅原始修复启发式消费
（\srcpath{src/scuc/scuc.cpp:build\_scuc\_primal\_repair\_seed}），
不进 MILP 约束\\
\texttt{PF} 变量块 & 分配 $n_lT$ 个连续变量但无任何约束/目标引用
（\srcpath{src/scuc/scuc.cpp:setup\_var\_index}）\\
负荷/风/光曲线缺行缺列 & 回退基准值/装机容量
（\srcpath{src/scuc/scuc.cpp:load\_at}、\srcpath{src/scuc/scuc.cpp:wind\_at}、
\srcpath{src/scuc/scuc.cpp:solar\_at}）\\
$B_{red}$ 不可逆 & \srcpath{src/scuc/scuc.cpp:compute\_ptdf} 取零逆矩阵，
PTDF 全 0，线路约束退化为无界行\\
母线/元件索引越界 & 装配期静默跳过（PTDF、潮流行、断面权重等）\\
硬编码：平衡母线 0 & \srcpath{src/scuc/scuc.cpp:compute\_ptdf}\\
硬编码：线路限额 $>10^5$ MW 不加约束 & \srcpath{src/scuc/scuc.cpp:build\_formulation}\\
硬编码：松弛变量上界 $10^6$、切负荷/过剩上界 $10^6$ & 同左\\
硬编码：$\mathrm{large\_uc}$ 阈值 $n_gT_c>800$ & 同左（关闭割平面族 T/H/R/P/Q）\\
硬编码：调节备用 10 分钟口径（$10\times$ MW/min） & 同左\\
硬编码：种子可行性容差 $10^{-5}$、数值零阈值 $10^{-9}$ & 同左、
\srcpath{src/scuc/scuc.cpp:scuc\_seed\_satisfies}\\
\fld{SOC}/\fld{soc\_init} $>1+10^{-8}$ 按 MWh 解释 & 同左
（初始 SOC 单位双口径）\\
\bottomrule
\end{longtable}
\end{footnotesize}
\end{center}


% numerical_validation.tex —— scuc 数值验证证据
% 规范：docs/modules/README.md 数值证据规范——只引用真实存在的测试目标、
% 基线文件与可复现命令；未执行的交叉验证不得写成已完成。

\section{数值验证}

\subsection{测试目标注册（CMake 事实）}

根 \texttt{CMakeLists.txt} 注册两个覆盖本模块的测试可执行文件（均要求
\texttt{MIPSOLVERS\_BUILD\_SCUC=ON}，默认开；运行时目录为 \texttt{tests/}）：

\begin{center}
\begin{footnotesize}
\begin{tabular}{@{}llll@{}}
\toprule
\textbf{目标} & \textbf{源文件} & \textbf{ctest 标签} & \textbf{说明}\\
\midrule
\texttt{test\_scuc\_module} & \srcpath{tests/test\_scuc\_module.cpp} &
\texttt{integration} & 模块级单元/集成测试\\
\texttt{test\_market\_simulation} & \srcpath{tests/test\_market\_simulation.cpp} &
\texttt{benchmark}（\texttt{TIMEOUT 600}） & 多求解器基准与物理一致性\\
\bottomrule
\end{tabular}
\end{footnotesize}
\end{center}

两者都直接编译 \srcpath{src/scuc/scuc.cpp}（后者还编译
\srcpath{src/scuc/case\_builder.cpp}）并链接 \texttt{mipsolvers} 与 Catch2。
复现命令（在已配置的构建目录下）：

\begin{verbatim}
ctest -R test_scuc_module --output-on-failure
ctest -R test_market_simulation --output-on-failure
# 或直接运行（Catch2 标签过滤）：
./tests/test_scuc_module "[scuc]"
./tests/test_market_simulation "[market][physics]"
\end{verbatim}

\subsection{覆盖矩阵：\texttt{test\_scuc\_module}}

\srcpath{tests/test\_scuc\_module.cpp} 以 2 母线、2 机组、3 时段、限额 150 MW
单支路的小问题（\srcpath{tests/test\_scuc\_module.cpp:make\_small\_problem\_json}，
\fld{mip\_gap}=0.005、时限 60 s）覆盖：

\begin{center}
\begin{footnotesize}
\begin{tabular}{@{}ll@{}}
\toprule
\textbf{TEST\_CASE（标签）} & \textbf{断言内容}\\
\midrule
JSON round-trip（\texttt{[scuc][unit]}） &
解析无异常；机组/支路/负荷数、\fld{num\_periods}、\fld{num\_buses}、
分段报价读回正确\\
small 2-bus problem converges（\texttt{[scuc][integration]}） &
SCUC 收敛、目标 $>0$、组合矩阵形状 $[2][3]$、切负荷 $|P^{shed}|<1$ MW、
时段 0 功率平衡（容差 5 MW）、能量/总成本 $>0$\\
SCED LP refines dispatch（\texttt{[scuc][integration]}） &
SCED 收敛且组合与 SCUC 一致（固定二进制口径，容差 0.01）\\
LMP values are physically consistent（\texttt{[scuc][integration]}） &
$0\le$ avg LMP $\le 300$ \$/MWh、min $\le$ avg $\le$ max、矩阵形状
$[n_b][T]$、高负荷时段 LMP 不低于低负荷时段（容差 0.1）\\
output JSON serialisation round-trips（\texttt{[scuc][unit]}） &
输出 JSON 可解析、含 \texttt{scuc}/\texttt{meta}/\texttt{lmp} 块且
收敛标志一致\\
\bottomrule
\end{tabular}
\end{footnotesize}
\end{center}

\subsection{覆盖矩阵：\texttt{test\_market\_simulation}}

\srcpath{tests/test\_market\_simulation.cpp} 基于 case\_builder 算例组织
（基准族对所有\textbf{当前构建可用}的 MILP 后端逐一运行，
\srcpath{tests/test\_market\_simulation.cpp:available\_milp\_solvers}；单后端
失败不致命、只记录，Auto 模式必须收敛）：

\begin{center}
\begin{footnotesize}
\begin{tabular}{@{}ll@{}}
\toprule
\textbf{TEST\_CASE（标签）} & \textbf{内容}\\
\midrule
3-bus 全求解器性能基准 T=6（\texttt{[market][benchmark]}） &
各后端测速对比；Auto 必须收敛\\
6-bus 风电+储能 全求解器基准 T=8（\texttt{[market][benchmark]}） &
含 \fld{use\_binary\_indicators} 之外全部储能路径（$M_2=50$）\\
IEEE 39-bus 快速基准 T=4 / 24h 风光 / 24h 全资源
（\texttt{[market][benchmark][ieee39]}） & 39 母线规模收敛与用时\\
备用约束验证（\texttt{[market][physics]}） &
$r^{spin}=0.10$、$r^{up}=r^{dn}=0.05$、$r^{neg}=0.10$、$R^{pfr}=10$ MW
下逐时段旋转备用 $\ge r^{spin}D_t-1$ MW\\
Gantt/出力/LMP 热图/SOC 时序（\texttt{[market][viz]}） &
ASCII 可视化 + 物理一致性抽查\\
IEEE 39-bus 24h 完整流程（\texttt{[market][viz][ieee39]}） &
三阶段全跑通 + JSON 导出\\
MIP Gap 收敛速度分析（\texttt{[market][benchmark]}） & 间隙-时间曲线记录\\
StrictHiGHS 39-bus root IPM + crossover 正确性
（\texttt{[market][strict\_highs][ipm\_root][ieee39]}） &
经 \srcpath{src/scuc/scuc.cpp:build\_scuc\_mip} 取裸模型直调
\texttt{solve\_milp\_bc}：IPM 根 + crossover 与单纯形根目标相对差 $<0.5\%$\\
StrictHiGHS 纯 IPM（无 crossover）（\texttt{...},\texttt{[pure\_ipm]}） &
记录无暖启动的节点 LP 代价；与 IPM+crossover 目标差 $<0.5\%$\\
\bottomrule
\end{tabular}
\end{footnotesize}
\end{center}

\subsection{CLI 冒烟用法}

构建后（可执行文件在构建输出目录；安装后经 \texttt{bin/}）：

\begin{verbatim}
# 生成算例 JSON
scuc_case_builder --case 3bus --T 6 --output case_3bus.json
scuc_case_builder --case ieee39 --T 24 --wind --solar --output case_39.json
# 求解（stderr 打印配置摘要与各阶段收敛状态）
scuc_solve case_3bus.json result_3bus.json --solver Auto --indent 2
scuc_solve case_39.json  result_39.json  --solver StrictHiGHS
# 仅 SCUC 阶段：--no-sced --no-lmp
\end{verbatim}

\texttt{scuc\_solve} 在 SCUC 未收敛时返回非零退出码（
\srcpath{src/scuc/main.cpp:main}），可直接接入调度脚本做失败告警；
\texttt{--solver} 覆盖 JSON 中的 \fld{solver} 字段。

\subsection{已知覆盖边界}

\begin{boundarynote}
\begin{itemize}
  \item 上述测试覆盖：JSON 往返、小规模收敛与功率平衡、SCED 组合一致性、
        LMP 数量级与单调性、备用需求、多后端基准、StrictHiGHS 根 LP 内核
        交叉校核（0.5\% 容差）；
  \item \textbf{未覆盖}：IEEE 118 母线规模无 \texttt{tests/} 注册测试
        （其求解器对比由 \srcpath{benchmark/bench\_01\_ieee118\_solvers.py}
        等 benchmark 脚本承载，不属于 ctest 套件）；无与外部参考求解器
        （如商业 MILP）在本仓库内的系统性目标值比对记录；N-1、多场景、
        交流潮流无实现亦无测试；断面（\fld{sections}）、机组群
        （\fld{generator\_groups}）、直流线、二进制储能模式
        （\fld{use\_binary\_indicators}）、热/温/冷启动字段无专门断言；
  \item 基准族对单后端失败\textbf{宽容}（仅记录），只有 Auto 模式的收敛是
        硬断言；阅读基准输出时须区分"求解器不可用"与"求解失败"；
  \item 容差口径：功率平衡 5 MW、LMP 单调性 0.1 \$/MWh、IPM/单纯形目标
        比对 0.5\%，均为小算例口径，不外推到大算例。
\end{itemize}
\end{boundarynote}


\section{接口与参数}

本章列出 \srcpath{src/scuc/scuc.cpp:scuc\_from\_json} 实际解析的全部 JSON 输入
字段与 \srcpath{src/scuc/scuc.cpp:scuc\_output\_to\_json} 实际写出的输出字段。
未列出即未消费；"默认值"指字段缺失时 C++ 结构体的初始化值。

\subsection{顶层结构}

\begin{paramtable}
\fld{config} & — & — & — & 见下表 & 求解配置块（可整体缺省）\\
\fld{num\_buses} & $n_b$ & 个 & $\ge 1$ & 自动推断 & 缺省时取所有元件 bus 最大值 +1\\
\fld{generators} & — & — & — & 空 & 机组数组，见 \ref{sec:par-gen} 节\\
\fld{branches} & — & — & — & 空 & 交流支路数组\\
\fld{loads} / \fld{wind} / \fld{solar} & — & — & — & 空 & 负荷/风电/光伏数组\\
\fld{storage} & — & — & — & 空 & 储能数组\\
\fld{dc\_lines} & — & — & — & 空 & 直流线路数组\\
\fld{generator\_groups} & — & — & — & 空 & 机组群数组\\
\fld{sections} & — & — & — & 空 & 监视断面数组\\
\fld{profiles} & — & — & — & 空 & 时序曲线块（load/wind/solar 三个矩阵）\\
\fld{initial\_status} & — & — & — & 空 & 初始状态块\\
\end{paramtable}

\subsection{\fld{config} 求解配置}

\begin{paramtable}
\fld{solver} & — & — & Auto/StrictHiGHS/ HiGHS/SCIP/ Native\-Branch\-And\-Cut & Auto & MILP 后端选择；大小写敏感\\
\fld{allow\_fallback} & — & — & true/false & true & 首选求解器失败时回退下一可用后端\\
\fld{num\_periods} & $T$ & 个 & 1--8760（CLI 校验） & 24 & 调度时段数\\
\fld{period\_length\_hr} & $\Delta t$ & h & $(0,24]$（CLI 校验） & 1.0 & 每时段小时数\\
\fld{n\_segments} & $K$ & 段 & $\ge 1$ & 3 & 每机组分段报价段数（取 $\max(1,\cdot)$）\\
\fld{mip\_gap} & $\varepsilon_{\mathrm{gap}}$ & — & $10^{-4}$--$10^{-2}$ & 0.001 & 相对 MIP 最优间隙容差\\
\fld{time\_limit\_sec} & — & s & — & 300 & 求解时限\\
\fld{verbose} & — & — & — & false & 后端日志开关\\
\fld{spinning\_reserve\_req} & $r^{\mathrm{spin}}$ & pu & 0--0.2 & 0.10 & 旋转备用需求（系统负荷比例）\\
\fld{regulation\_up\_req} & $r^{\mathrm{up}}$ & pu & 0--0.1 & 0.05 & 上调调节备用比例\\
\fld{regulation\_down\_req} & $r^{\mathrm{dn}}$ & pu & 0--0.1 & 0.05 & 下调调节备用比例\\
\fld{neg\_reserve\_req} & $r^{\mathrm{neg}}$ & pu & 0--0.2 & 0.0 & 负备用需求比例；0 = 不加该约束\\
\fld{pfr\_reserve\_req\_mw} & $R^{\mathrm{pfr}}$ & MW & 0--数百 & 0.0 & 一次调频备用需求；0 = 不加该约束\\
\fld{voll} & $c^{\mathrm{shed}}$ & \$/MWh & $10^3$--$10^5$ & 10000 & 切负荷惩罚（失负荷价值）\\
\fld{vocc} & $c^{\mathrm{curt}}$ & \$/MWh & $10^2$--$10^3$ & 1000 & 系统级发电过剩惩罚\\
\fld{renewable\_min\_output\_coeff} & $\alpha_n$ & pu & 0--1 & 0.0 & 可再生最低出力系数（预测比例）\\
\fld{M1\_line\_slack\_penalty} & $M_1$ & \$/MW & $10^4$--$10^6$ & $10^5$ & 线路/断面潮流松弛罚系数\\
\fld{M2\_renewable\_curtail\_penalty} & $M_2$ & \$/MWh & 0--$10^2$ & 0.0 & 弃风/弃光惩罚；$>0$ 才启用弃电变量\\
\fld{wheeling\_fee\_per\_mwh} & $c^{\mathrm{wf}}$ & \$/MWh & 0--数十 & 0.0 & 输电费（加在全部机组出力上）；0 = 停用\\
\fld{enable\_market\_cuts} & — & — & — & true & 预构型市场割平面开关（第 4 章割平面族）\\
\fld{enable\_primal\_repair} & — & — & — & true & UC 承诺种子 + 连续修复启发式开关\\
\fld{solve\_sced} & — & — & — & true & SCUC 后是否执行 SCED 阶段\\
\fld{solve\_lmp} & — & — & — & true & 是否执行 LMP 阶段\\
\fld{lmp\_delta} & $\delta$ & pu & 0.05--0.2 & 0.10 & LMP 重调度邻域因子\\
\end{paramtable}

\subsection{\fld{generators} 机组字段}
\label{sec:par-gen}

\begin{paramtable}
\fld{name} & — & — & — & 空串 & 机组名（仅标识用）\\
\fld{bus} & $b_g$ & — & $[0,n_b)$ & 0 & 接入母线（0 起始）\\
\fld{pmin} / \fld{pmax} & $\underline P_g,\bar P_g$ & MW & — & 0 & 最小/最大技术出力\\
\fld{ramp\_up\_mw\_min} & $R^{+}_g$ & MW/min & — & 0 & 上爬坡率（装配时乘 $60\Delta t$）\\
\fld{ramp\_dn\_mw\_min} & $R^{-}_g$ & MW/min & — & 0 & 下爬坡率\\
\fld{min\_up\_time\_hr} & $T^{U}_g$ & h & — & 0 & 最小开机时间（装配时向上取整）\\
\fld{min\_dn\_time\_hr} & $T^{D}_g$ & h & — & 0 & 最小停机时间\\
\fld{must\_run} & — & — & — & false & 必开机组（组合变量固定为 1）\\
\fld{max\_startups} & — & 次 & — & 0 & 日内启动次数上限；0 = 不限\\
\fld{max\_shutdowns} & — & 次 & — & 0 & 日内停机次数上限；0 = 不限\\
\fld{bid\_segments} & $(q_k,c_k)$ & MW, \$/MWh & — & 空 & 分段报价（Pmin 之上逐段容量与边际价）\\
\fld{startup\_cost} & $C^{U}_g$ & \$/次 & — & 0 & 启动成本（当前实现对所有启动状态均用此值）\\
\fld{startup\_cost\_warm} / \fld{startup\_cost\_cold} & — & \$/次 & — & 0 & \textbf{已解析未消费}：温/冷启动成本未进入模型\\
\fld{hot\_start\_threshold\_hr} / \fld{warm\_start\_threshold\_hr} & — & h & — & 4 / 8 & \textbf{已解析未消费}：热/温启动判定时长\\
\fld{no\_load\_cost} & $C^{NL}_g$ & \$/h & — & 0 & 空载成本（目标中乘 $1/\mathrm{iph}$）\\
\fld{spinning\_reserve\_price} & $c^{\mathrm{spin}}_g$ & \$/MWh & — & 0 & 旋转备用报价\\
\fld{regulation\_up\_price} / \fld{regulation\_down\_price} & $c^{\mathrm{up}}_g,c^{\mathrm{dn}}_g$ & \$/MWh & — & 0 & 调节备用报价\\
\fld{pfr\_alpha} & $\alpha^{\mathrm{pfr}}_g$ & pu & 0--0.1 & 0 & 一次调频系数；0 = 不参与\\
\fld{ud\_periods} / \fld{dd\_periods} & — & 时段 & — & 0 & \textbf{已解析未消费}：启动/停机过程时长未建模\\
\fld{group\_id} & — & — & — & $-1$ & \textbf{已解析未消费}：群约束以 \fld{generator\_groups} 为准\\
\end{paramtable}

\subsection{网络、负荷与可再生字段}

\begin{paramtable}
\fld{branches[].from}/\fld{branches[].to} & — & — & $[0,n_b)$ & 0 & 支路端点（0 起始）\\
\fld{branches[].reactance} & $x_l$ & pu & $10^{-3}$--1 & 0.05 & 电抗；$|x|<10^{-6}$ 时按 $10^{-3}$ 兜底\\
\fld{branches[].rating\_mw} & $\bar F_l$ & MW & — & $10^6$ & 热稳定限额；$>10^5$ 的支路不加潮流约束\\
\fld{branches[].in\_service} & — & — & — & true & 退出运行的支路不进 Bbus 与潮流约束\\
\fld{loads[].bus} / \fld{loads[].p\_mw} & $D_d$ & MW & — & 0 / 0 & 基准负荷；实际负荷 = 基准 × 负荷曲线 pu\\
\fld{wind[].bus} / \fld{wind[].pmax} & $\bar W_w$ & MW & — & 0 / 0 & 风电场；曲线缺省时预测取 pmax\\
\fld{solar[].bus} / \fld{solar[].pmax} & $\bar S_s$ & MW & — & 0 / 0 & 光伏场；同上\\
\fld{dc\_lines[].from}/\fld{dc\_lines[].to} & — & — & — & 0 & 直流线注入/受端母线\\
\fld{dc\_lines[].pmin}/\fld{dc\_lines[].pmax} & — & MW & — & $\mp 10^6$ & 直流功率上下限（可负，表示反向送电）\\
\fld{dc\_lines[].ramp\_up}/\fld{dc\_lines[].ramp\_dn} & — & MW/时段 & — & $10^6$ & 相邻时段功率变化限（首时段相对 0）\\
\end{paramtable}

\subsection{\fld{storage} 储能字段}

\begin{paramtable}
\fld{bus} & — & — & — & 0 & 接入母线\\
\fld{pmax\_charge} / \fld{pmax\_discharge} & $\bar P^{ch},\bar P^{dis}$ & MW & — & 0 & 充/放功率上限\\
\fld{pmin\_charge} / \fld{pmin\_discharge} & $\underline P^{ch},\underline P^{dis}$ & MW & — & 0 & 充/放最小功率（仅二进制模式下生效）\\
\fld{energy\_capacity\_mwh} & $\bar E$ & MWh & — & 0 & 能量容量\\
\fld{efficiency} & $\eta$ & — & 0.8--0.95 & 0.9 & 往返效率（钳到 $[0.1,1]$）\\
\fld{eta\_charge} / \fld{eta\_discharge} & $\eta^{ch},\eta^{dis}$ & — & — & $-1$ & 充/放分效率；$<0$ 时取 $\sqrt{\eta}$\\
\fld{soc\_init} & $E_0$ & pu 或 MWh & 0--1 & 0.5 & 初始 SOC；$>1$ 时按 MWh 解释\\
\fld{soc\_min} & $\underline E$ & pu & 0--1 & $-1$ & SOC 下限；$<0$ 时取 0.10\\
\fld{soc\_final} & $E^{fin}$ & pu & 0--1 & $-1$ & 末时段 SOC 要求；$<0$ 时取初始 SOC\\
\fld{charge\_bid\_price} / \fld{discharge\_bid\_price} & $\lambda^{ch},\lambda^{dis}$ & \$/MWh & — & 0 & 充/放电报价（目标中乘 $\Delta t$；绝对值 $>\varepsilon$ 才写入）\\
\fld{cycle\_limit} & $N^{cycle}$ & 次 & 0--2 & 0 & 日等效全循环上限；0 = 不限\\
\fld{use\_binary\_indicators} & — & — & — & false & 启用 $\xi^+/\xi^-$ 二进制充放模式（严格互斥）\\
\end{paramtable}

\subsection{\fld{generator\_groups}、\fld{sections}、\fld{profiles}、\fld{initial\_status}}

\begin{paramtable}
\fld{generator\_groups[].gen\_indices} & — & — & — & 空 & 成员机组索引（0 起始）\\
\fld{generator\_groups[].pmin\_t}/\fld{pmax\_t} & — & MW & — & 空 & 群出力逐时段上下限；长度 1 时全时段共用；空 = 不约束\\
\fld{generator\_groups[].emin}/\fld{emax} & — & MWh & — & 0 & 群全时段电量上下限；0 = 不约束\\
\fld{sections[].line\_weights} & $w_l$ & — & — & 空 & 成员（线路号， 权重）；接受 \texttt{[line,weight]} 或 \texttt{\{line,weight\}} 两种 JSON 形态\\
\fld{sections[].rating\_fwd\_mw}/\fld{rating\_rev\_mw} & — & MW & — & $10^6$ & 断面正/反向限额\\
\fld{profiles.load} & — & pu & 0.5--1.2 & 空 & $\mathrm{load}[d][t]$：负荷 $d$ 在时段 $t$ 的 pu 倍率；缺行/缺列回退基准值\\
\fld{profiles.wind} / \fld{profiles.solar} & — & MW & — & 空 & 风/光\textbf{绝对}预测出力；缺省时回退 pmax\\
\fld{initial\_status.commitment} & $u_{g,0}$ & 0/1 & — & 空 & 时段 0 之前在/离网状态\\
\fld{initial\_status.dispatch} & $P_{g,0}$ & MW & — & 空 & 时段 0 之前出力（首时段爬坡基准）\\
\fld{initial\_status.storage\_soc} & — & pu/MWh & — & 空 & 储能初始 SOC（覆盖 \fld{soc\_init}）\\
\fld{initial\_status.time\_in\_state} & — & h & — & 空 & 当前状态已持续时间（正=开机，负=停机）；\textbf{仅原始修复启发式消费}，不进 MILP 约束\\
\end{paramtable}

\subsection{输出 JSON 字段}

由 \srcpath{src/scuc/scuc.cpp:scuc\_output\_to\_json} 与匿名命名空间内的
\srcpath{src/scuc/scuc.cpp:solve\_result\_to\_json} 写出。注意：
\texttt{SCUCSolveResult} 中的 \texttt{segment\_dispatch}、
\texttt{wind\_curtailment}、\texttt{solar\_curtailment}、\texttt{section\_flows}
四个矩阵\textbf{不进入} JSON 输出（仅 C++ 层可见）。

\begin{paramtable}
\fld{meta.solver} 等 & — & — & — & — & 输入元信息回显（solver、num\_periods、num\_buses、num\_generators、num\_branches）\\
\fld{scuc.converged} / \fld{sced.converged} & — & — & — & — & 阶段收敛标志；false 时矩阵为空、标量为 0\\
\fld{scuc.objective} & $z$ & \$ & — & — & 求解器返回的目标值（含全部罚项）\\
\fld{scuc.solver\_name} & — & — & — & — & 实际使用的后端名；SCED 阶段追加 \texttt{" (SCED-LP)"}\\
\fld{scuc.solve\_time\_sec} & — & s & — & — & 阶段墙钟用时\\
\fld{scuc.mip\_gap} & — & — & — & — & MIP 间隙（SCED/LMP 为 LP，无此口径）\\
\fld{scuc.n\_cuts\_added} & — & 行 & — & — & 预构型割平面行数\\
\fld{scuc.commitment}/\fld{startup}/\fld{shutdown} & $u,y,z$ & 0/1 & $[n_g][T_c]$ & — & 组合/启动/停机矩阵\\
\fld{scuc.dispatch} & $P_{g,t}$ & MW & $[n_g][T]$ & — & 机组出力\\
\fld{scuc.spinning\_reserve} 等 & $R$ & MW & $[n_g][T]$ & — & 三类备用（spinning/regulation\_up/regulation\_down）\\
\fld{scuc.wind\_generation}/\fld{solar\_generation} & — & MW & $[n_w][T]$ 等 & — & 可再生出力\\
\fld{scuc.storage\_charging}/\fld{storage\_discharging}/\fld{storage\_soc} & — & MW/MWh & $[n_s][T]$ & — & 储能充放功率与 SOC\\
\fld{scuc.line\_flows} & $F_{l,t}$ & MW & $[n_l][T]$ & — & 由 PTDF × 节点注入重算的支路潮流\\
\fld{scuc.load\_shedding}/\fld{gen\_curtailment} & — & MW & $[T]$ & — & 系统级切负荷/过剩\\
\fld{scuc.cost.*} & — & \$ & — & — & 成本分解：energy/startup/no\_load/reserve/penalty/total（口径见第 4 章）\\
\fld{lmp.nodal}/\fld{energy}/\fld{congestion} & $\lambda$ & \$/MWh & $[n_b][T]$ & — & 节点电价及能量/阻塞分量\\
\fld{lmp.avg\_lmp}/\fld{max\_lmp}/\fld{min\_lmp} & — & \$/MWh & — & — & 时空汇总\\
\end{paramtable}

\section{验证与边界局限}

\subsection{验证状态}

模块的数值验证证据集中在第 5 章：JSON 往返、2 母线 3 时段收敛与功率平衡、
SCED 组合一致性、LMP 物理合理性、备用约束、多求解器基准与 StrictHiGHS
根节点 IPM 正确性均有注册测试覆盖；IEEE 118 母线规模有 benchmark 族算例
（\srcpath{benchmark/}，见第 5 章边界声明）。

\subsection{边界与局限}

\begin{boundarynote}
\begin{itemize}
  \item \textbf{网络模型}：直流潮流近似（无损、电压幅值恒定、小相角），平衡母线
        固定为 0 号；额定 $>10^5$ MW 的支路不加潮流约束（视为不阻塞）；无 N-1
        预想事故集、无交流潮流校验；
  \item \textbf{时间耦合}：最小启停约束在时界起点截断（$\tau\ge 0$），初始
        在/离网历史（\fld{time\_in\_state}）不进 MILP 约束，仅原始修复启发式
        使用；
  \item \textbf{启动成本}：热/温/冷三态字段已解析但未消费，所有启动均按
        \fld{startup\_cost} 计费；启停过程出力曲线（\fld{ud\_periods}/
        \fld{dd\_periods}）未建模；
  \item \textbf{LMP}：取自固定组合 LP 的对偶，含 $\delta$ 邻域限幅的再调度
        口径；MILP 本身不返回有意义的行对偶，不要从 SCUC 阶段直接取价；
  \item \textbf{规模}：割平面族 T/H/R/P/Q 在 $n_g \times T_c > 800$ 时自动关闭
        （行开销大于收紧收益，见第 4 章）；PTDF 为稠密 $n_l \times n_b$ 矩阵，
        超大系统内存随其平方增长；
  \item \textbf{异常}：输入 JSON 语法/类型错误抛 nlohmann 异常；求解失败沿
        \texttt{converged=false} 如实上报，输出矩阵为空、标量为 0，不得以零值
        误读为可行解。
\end{itemize}
\end{boundarynote}

\section{跨模块工业级技术评价基线}

本章规定本手册所述模块进入工程算例、批量研究或生产服务前必须共同满足的评价基线。
具体模型公式、算法步骤、选项和终止判据以正文相应章节为准；本章不扩大源码已经实现的范围。

\subsection{输入、输出、边界条件与数据结构审计}
输入审计应逐项检查问题维度一致性、系数有限性、上下界合法性（含 $\ell\le u$）、
变量类型标记（连续/整数/二进制）、锥成员归属与选项枚举值；非法输入必须在进入求解内核前由
\srcpath{src/engine/util/problem\_validation.cpp} 一类入口拒绝并返回结构化错误。
输出审计必须记录求解状态（最优/不可行/无界/时限/数值异常且相互区分）、目标值、
原始与对偶残差、间隙、后端名称、容差、迭代/节点计数和告警。
空问题、零行零列、退化约束、固定变量、重复索引与不可用可选后端均属于边界测试集。

\subsection{工程假设与数值稳定性分析}
模型使用的凸性、光滑性、约束规范、中心路径存在性、基非奇异性等假设必须与所求解的问题类匹配。
数值评价至少包含数据缩放（scaling）、原始/对偶不可行度、互补间隙、KKT 残差、
基或因子分解的条件数/枢轴质量、步长/阻尼、迭代停滞与容差复现性。
若采用正则化、扰动、惯性校正、回退后端或近似（如容差放宽），应同时报告触发条件、参数值、
对主指标的敏感性以及是否改变结果口径；不得用``已返回结果''替代最优性或收敛性证明。

\subsection{复杂度与资源分析}
令 $m$、$n$、$n_z$ 分别表示约束数、变量数与矩阵非零元数。
报告应分别给出建模、预求解、求解、后处理的时间与内存。
单纯形单次迭代的稠密代价为 $O(m^2)$ 量级，实际取决于基因子更新与定价策略；
内点法每步代价由 KKT 稀疏因子分解主导，应报告结构非零数、填充率和因子分解峰值内存；
MILP 不给出虚假的多项式最坏界，而报告节点数、间隙、时限和实测扩展曲线。
复杂度结论必须基于固定硬件、固定后端、固定构建配置和可复现算例族。

\subsection{异常工况处理}
不可行、无界、数值失败、时限、内存不足和后端缺失必须相互区分并沿 API 如实上报；
不得把数值失败静默改写为不可行，不得把次优解标为全局最优，不得用零值或 NaN 伪装未知量。
异常路径必须保证问题对象可恢复、可重试；回退（fallback）到替代后端时必须在结果或日志中
声明实际使用的求解器。

\subsection{与其他模块的接口关系}
接口链路遵循``AML/文件输入 $\to$ 问题校验 $\to$ 预求解 $\to$ 求解内核或外部后端
$\to$ 后处理/解回写''。跨模块传递时保留变量/约束的稳定身份、索引空间、符号约定与单位；
消费 SCUC 等上层应用结果前，先检查其求解状态与容差口径。Python 绑定层只做语义保持适配，
不重新解释求解结果。

\subsection{测试与验证建议}
最低验证矩阵包括：解析可解小例、标准算例集（如 NETLIB、MIPLIB 2017）、边界/异常输入、
算法或后端交叉校核、序列化往返、重复运行确定性和规模扩展测试。
数值算法还应加入残差独立重算；近似路径应与高保真路径比较并给出误差分布，而非只报告平均值。
回归报告必须列明提交哈希、构建配置、算例版本、容差、通过/失败/跳过数及未验证范围。

\subsection{工业级评估指标}
\begin{itemize}
  \item \textbf{正确性}：KKT 残差、约束最大违反、互补间隙、解析/基准目标误差、不可行证书正确性；
  \item \textbf{鲁棒性}：标准算例集收敛率、不可行/无界识别率、数值失败率、容差敏感性；
  \item \textbf{性能}：端到端时延、迭代/节点计数、峰值内存、并行效率与规模增长斜率；
  \item \textbf{可追溯性}：选项快照、后端与版本记录、容差口径、告警可定位率；
  \item \textbf{工程有效性}：与参考求解器或历史基线的目标值/时间偏差，及其对业务指标的影响。
\end{itemize}
指标应预先给出验收阈值和失败处置；未达到阈值时只能标记为实验性或受限适用。

\subsection{适用范围、局限性与后续改进}
本手册只评价当前源码覆盖的问题类、后端和选项，不对未实现的问题类、未安装求解器或未执行的
外部验证作保证。后续改进应按``缺口 $\to$ 可量化指标 $\to$ 源码/测试变更 $\to$ 独立验证''闭环，
优先消除会造成错误结论的语义缺口，其次提升数值鲁棒性与性能，最后扩展问题类范围。


\end{document}
